% Generated semantic source for AI Almanach v3.
% Edit v3 directly after initial generation; do not overwrite
% editorial changes by rerunning the converter casually.

% ==== Natural Language Processing and Generative AI =============================================
\chapter{Natural Language Processing and Generative AI}

\chapterlead{How can language systems be evaluated rather than merely admired?}{Tokenise}{Represent}{Generate}{Evaluate}{nlp-generative-ai}

% ==== Section 1: Natural Language Processing Foundations =====================================
\label{ch:nlp-voice-assistants}

\begin{learningobjectives}
After working through this chapter you should be able to:
\begin{itemize}
  \item run the byte-pair-encoding merge loop by hand and predict how an unseen
        word will be split;
  \item explain why token count, not word count, is what a context window
        measures --- and estimate both;
  \item compute perplexity from a sequence of probabilities and relate it to
        cross-entropy;
  \item compute a BLEU score, including the brevity penalty, and state what it
        cannot see;
  \item estimate an $n$-gram probability with and without smoothing, and say
        what smoothing costs;
  \item describe what a word embedding does and does not encode;
  \item size a retrieval-augmented pipeline: chunk count, retrieved context,
        and what is left for the answer;
  \item name the two separate evaluations a RAG system needs;
  \item explain why a fluent answer is not evidence of a correct one.
\end{itemize}
\end{learningobjectives}

\begin{plainlanguage}
A language model does not read words. It reads numbers. The whole first half of
this chapter is about how text becomes numbers --- splitting it into tokens,
mapping tokens to IDs, and turning IDs into vectors that place similar meanings
near each other. The second half is about what you can build once that
translation exists: translation, search, summaries, assistants. Almost every
surprising behaviour of a language model, good or bad, traces back to a
decision made in that first translation step.
\end{plainlanguage}

\begin{engineernote}
This chapter sits between \cref{ch:deep-learning}, which builds the Transformer,
and \cref{ch:practical-llm-engineering}, which takes a real checkpoint apart.
Read it as the bridge: here the mechanism becomes a language system, and the
failures stop being numerical and start being linguistic and social.
\end{engineernote}

\section{Natural Language Processing Foundations}

\subsection{From Pixels to Words: Why Language Is Harder}

\begin{v3topic}{}
It is worth being explicit about why language needed different machinery from
the image problems of \cref{ch:computer-vision}, because each difference forces
a specific architectural choice --- and together they amount to a derivation of
the architecture that answered
them\textsuperscript{\cite{sutskeverSequenceSequenceLearning2014},
\cite{bahdanauNeuralMachineTranslation2015},
\cite{vaswaniAttentionAllYou2023}}.

\begin{table}[htbp]
\centering
\caption{What makes language structurally harder than a fixed-size image
classification problem, and the mechanism each difficulty forced into
existence.}
\label{tab:nlp-why-language-is-harder}
\begin{tblr}{
  width=\textwidth,
  colspec={Q[l,0.20\textwidth] X[l] X[l]},
  row{1}={font=\bfseries,bg=AlmanachMist},
  hlines,vlines}
Property & Images & Language, and what it forces \\
Input size & Fixed: every MNIST digit is $784$ pixels &
  Variable: five words or five thousand. Forces a mechanism that works at any
  length \\
Structure & A 2D grid; neighbours are related by construction &
  No grid. ``The dog bit the man'' and ``the man bit the dog'' share every word
  and differ in meaning, so \emph{order} must be encoded explicitly \\
Dependency range & Local; convolution's small kernels suffice &
  Arbitrary. In ``the cat that the dog chased ran away'', \emph{ran} belongs to
  \emph{cat}, not the nearer \emph{dog}. Forces a mechanism with constant path
  length between any two positions \\
Output & One of ten labels &
  A sequence, generated one token at a time from a vocabulary of $30{,}000$ or
  more. Forces autoregressive decoding \\
\end{tblr}
\end{table}

\begin{engineernote}
Read the right-hand column as a derivation of the Transformer rather than a
list of complaints. Variable length rules out a fixed input layer; the absence
of a grid rules out convolution as the primary mixer and requires positional
encoding; arbitrary dependency range is precisely what
self-attention's $O(1)$ path length provides
(\cref{tab:dl-seq-complexity}); and generation rather than classification is
what makes the causal mask necessary. Nothing about the architecture is
arbitrary --- each piece answers one row of this table.
\end{engineernote}
\end{v3topic}

% ---- 12.1.1 What is NLP? -------------------------------------------------------------------
\subsection{What is NLP?}

\begin{v3topic}{Natural Language Processing}
    \textbf{Natural Language Processing (NLP)} is the field at the intersection of computer science and linguistics concerned with building machines that can understand, manipulate, and generate human language\textsuperscript{\cite[][p.~3]{dhamaniIntroductionGenerativeAI2024}}.
    The primary task drivers are:
    \begin{itemize}
        \item \textbf{Understanding}: extract meaning from text or speech
        \item \textbf{Generation}: produce grammatically correct, coherent text
        \item \textbf{Translation}: map between languages
        \item \textbf{Dialogue}: maintain coherent multi-turn conversation
    \end{itemize}
    NLP is positioned within AI as a subfield of machine learning, itself enabled by deep learning.
    Large language models (LLMs) are the current dominant NLP technology.

\end{v3topic}
\begin{v3topic}{Core Challenges in NLP}
    Human language is inherently ambiguous and context-dependent\textsuperscript{\cite[][p.~9]{dhamaniIntroductionGenerativeAI2024}}:
    \begin{itemize}
        \item \textbf{Lexical ambiguity}: ``bank'' can mean financial institution or river bank
        \item \textbf{Syntactic ambiguity}: ``I saw the man with the telescope'' --- who has the telescope?
        \item \textbf{Co-reference}: ``Julia said she arrived'' --- who is ``she''?
        \item \textbf{Pragmatics}: ``Can you pass the salt?'' is a request, not a yes/no question
        \item \textbf{World knowledge}: ``The trophy did not fit because it was too big'' --- what is ``it''?
    \end{itemize}
    Humans resolve these effortlessly through context; teaching machines to do so remains central to NLP research.

\end{v3topic}
\begin{v3topic}{A Brief History of NLP}
    NLP development spans three eras\textsuperscript{\cite[][pp.~3--9]{dhamaniIntroductionGenerativeAI2024}}:

    \begin{tblr}{width=\linewidth,
        colspec = {X[1.2]X[3.5]},
        hlines, vlines,
        row{1} = {font=\bfseries\small, bg=LimeGreen!30},
        rows = {font=\small},
    }
        Period & Milestone \\
        1940s--1960s & Machine translation as first NLP goal; ELIZA chatbot (Weizenbaum, 1966) uses rule-based pattern matching \\
        1970s--1980s & Symbolic / rule-based NLP; ``AI winter'' limits progress \\
        1990s & Statistical methods dominate; models learn from data rather than hand-coded rules \\
        2006 & Google Translate --- first commercially successful NLP product \\
        2013 & Word2Vec (Google) --- dense word embeddings trained from text \\
        2017 & Transformer architecture (Vaswani et al.) --- foundation of modern NLP \\
        2018--2019 & GPT-1 (OpenAI) and BERT (Google) --- pre-trained large language models \\
        2020--2023 & GPT-3 (175B), ChatGPT, GPT-4 --- conversational AI at scale \\
    \end{tblr}
\captionof{table}{The successive paradigms of natural language processing, each of which solved the previous one's blocking limitation.}
\label{tab:nlp-history}
\end{v3topic}

% ---- 12.1.2 Linguistic Building Blocks -----------------------------------------------------
\subsection{Linguistic Building Blocks}

\begin{v3topic}{Levels of Linguistic Analysis}
    Language is analysed at multiple levels\textsuperscript{\cite[][pp.~3--4]{dhamaniIntroductionGenerativeAI2024}}.
    Each level provides different information for NLP tasks:

    \begin{tblr}{width=\linewidth,
        colspec = {Q[2.5cm]X[2]X[4]},
        hlines, vlines,
        row{1} = {font=\bfseries\small, bg=LimeGreen!30},
        rows = {font=\small},
    }
        Level & Studies & Example \\
        \textbf{Morphology} & Word forms & \emph{run} $\to$ \emph{ran}, \emph{running}; \emph{happy} $\to$ \emph{unhappiness} \\
        \textbf{Syntax} & Sentence structure & Parse tree; subject-verb-object order; dependency relations \\
        \textbf{Semantics} & Meaning & Word sense: ``bank'' (finance vs.\ river); compositionality \\
        \textbf{Pragmatics} & Language in context & Speaker intent; indirect requests; discourse coherence \\
        \textbf{Phonology} & Sound patterns & Stress, intonation, phoneme inventories (for speech) \\
    \end{tblr}
\captionof{table}{Levels of linguistic analysis. Modern end-to-end models do not represent these explicitly, but the levels remain the vocabulary for describing where a system fails.}
\label{tab:nlp-linguistic-levels}

    NLP systems typically must handle all levels: a machine translation system requires morphological analysis, syntactic parsing, semantic understanding, and pragmatic coherence.
\end{v3topic}

% ---- 12.1.3 Tokenization -------------------------------------------------------------------
\subsection{Tokenization}

\begin{v3topic}{Tokens: The Atoms of Language}
    Neural language models cannot process raw text directly; they require numeric inputs.
    A \textbf{token} is the smallest unit of text a model processes --- the ``atom'' of language\textsuperscript{\cite[][p.~15]{farrisHowLargeLanguage2025}}.
    \textbf{Tokenization} converts a text string into a sequence of integer token IDs:
    \[
        \text{``Hello World!''} \xrightarrow{\text{tokenize}} [4944,\ 3574,\ 0]
    \]
    Tokens are not necessarily words --- they are subword units that represent the \emph{minimum meaningful semantic unit} that the tokenizer has learned from training data.
    GPT-3 uses a vocabulary of 50{,}257 tokens; most common English words map to a single token, while rare words are decomposed into several subword tokens.

\end{v3topic}
\begin{v3topic}{The Tokenization Pipeline}
    The tokenization process follows four steps\textsuperscript{\cite[][p.~17]{farrisHowLargeLanguage2025}}:
    \begin{enumerate}
        \item \textbf{Input}: receive text string from user or document source
        \item \textbf{Normalization}: transform string --- convert to lowercase, handle Unicode (homoglyphs), remove or preserve punctuation
        \item \textbf{Segmentation}: split normalised string into subword tokens
        \item \textbf{Output}: map each token to its unique integer ID in the vocabulary
    \end{enumerate}
    The vocabulary is fixed after training: any token not seen during training must be decomposed from known subwords (the \textbf{out-of-vocabulary} problem).
    Normalization choices create a trade-off: aggressive lowercasing reduces vocabulary size but loses information (e.g.,\ \emph{Bill} vs.\ \emph{bill}).

\end{v3topic}
\begin{v3topic}{Byte Pair Encoding (BPE)}
    \textbf{BPE}\textsuperscript{\cite{sennrichNeuralMachineTranslation2016}} is the dominant subword tokenization algorithm, used in GPT (via \texttt{tiktoken}), RoBERTa, and others\textsuperscript{\cite[][pp.~20--22]{farrisHowLargeLanguage2025}}:
    \begin{enumerate}
        \item Initialise vocabulary with individual characters $\{a, b, \ldots, z\}$
        \item Find the most frequently co-occurring adjacent byte pair in the corpus
        \item Merge that pair into a new subword token; add to vocabulary
        \item Repeat until vocabulary reaches target size
    \end{enumerate}
    Result: common words become single tokens; rare words decompose.
    \emph{loquacious} $\to$ \texttt{lo} + \texttt{qu} + \texttt{acious}.
    \emph{schoolhouse} $\to$ \texttt{school} + \texttt{house}.
    Alternative algorithms: \textbf{WordPiece} (BERT, Google) and \textbf{SentencePiece} (multilingual models, preserves whitespace).

\begin{figure}[htbp]
\centering
\begin{tikzpicture}[
  bx/.style={draw=AlmanachTeal,fill=AlmanachMist,rounded corners=1mm,
    minimum height=8mm,font=\ttfamily\scriptsize,inner xsep=3pt},
  vec/.style={draw=AlmanachOchre,fill=AlmanachOchre!18,rounded corners=0.6mm,
    minimum height=8mm,minimum width=7mm,font=\scriptsize},
  arr/.style={-{Stealth[length=1.8mm]},thick,draw=AlmanachInk!55},
  lbl/.style={font=\sffamily\scriptsize,color=AlmanachInk!70}]

  \node[bx,minimum width=34mm] (txt) at (0,0) {"unhappiness"};
  \node[lbl,below=1.5mm of txt] {raw string};

  \node[bx] (t1) at (33mm,4.5mm) {un};
  \node[bx] (t2) at (33mm,-4.5mm) {happiness};
  \node[lbl] at (33mm,-13mm) {subword tokens};

  \node[bx] (i1) at (56mm,4.5mm) {392};
  \node[bx] (i2) at (56mm,-4.5mm) {21832};
  \node[lbl] at (56mm,-13mm) {integer IDs};

  \foreach \r in {0,1,2}{
    \node[vec] at (78mm+\r*8mm,4.5mm) {};
    \node[vec] at (78mm+\r*8mm,-4.5mm) {};}
  \node[lbl] at (94mm,-13mm) {embedding vectors};
  \node[lbl,anchor=west] at (102mm,4.5mm) {$\cdots$};
  \node[lbl,anchor=west] at (102mm,-4.5mm) {$\cdots$};

  \draw[arr] (txt.east) -- ++(4mm,0) |- (t1.west);
  \draw[arr] (txt.east) -- ++(4mm,0) |- (t2.west);
  \draw[arr] (t1) -- (i1);
  \draw[arr] (t2) -- (i2);
  \draw[arr] (i1.east) -- ++(4mm,0);
  \draw[arr] (i2.east) -- ++(4mm,0);
\end{tikzpicture}
\caption{How text becomes something a network can multiply. The model never
sees letters --- only the IDs in the third column, and then the learned vectors
in the fourth. Everything the model can and cannot do about spelling is decided
at the second arrow.}
\label{fig:nlp-tokenization-pipeline}
\end{figure}

\begin{workedexample}[Three BPE merges, by hand]
Train a tokenizer on a five-word corpus with these frequencies:
\texttt{hug}\,(10), \texttt{pug}\,(5), \texttt{hugs}\,(5), \texttt{pun}\,(12),
\texttt{bun}\,(4).

\emph{Start.} Split every word into characters. The base vocabulary is
$\{$\texttt{b}, \texttt{g}, \texttt{h}, \texttt{n}, \texttt{p}, \texttt{s},
\texttt{u}$\}$.

\emph{Round 1 --- count adjacent pairs.}
\[
\begin{aligned}
  \texttt{h}\,\texttt{u} &= 10+5 = 15, &
  \texttt{u}\,\texttt{g} &= 10+5+5 = \mathbf{20}, &
  \texttt{g}\,\texttt{s} &= 5,\\
  \texttt{p}\,\texttt{u} &= 5+12 = 17, &
  \texttt{u}\,\texttt{n} &= 12+4 = 16, &
  \texttt{b}\,\texttt{u} &= 4.
\end{aligned}
\]
The winner is \texttt{ug} at $20$. Merge it and add it to the vocabulary.

\emph{Round 2.} Recount with \texttt{ug} now a single symbol:
$\texttt{h}\,\texttt{ug} = 15$, $\texttt{p}\,\texttt{ug} = 5$,
$\texttt{ug}\,\texttt{s} = 5$, $\texttt{p}\,\texttt{u} = 12$,
$\texttt{u}\,\texttt{n} = \mathbf{16}$, $\texttt{b}\,\texttt{u} = 4$.
Merge \texttt{un}.

\emph{Round 3.} Now $\texttt{h}\,\texttt{ug} = \mathbf{15}$ is the largest
($\texttt{p}\,\texttt{un} = 12$ is next). Merge \texttt{hug}.

\emph{Result.} Vocabulary $=$ base characters $+\ \{$\texttt{ug},
\texttt{un}, \texttt{hug}$\}$.

\emph{Now tokenize new words with it:}
\[
\begin{aligned}
  \texttt{hugs} &\rightarrow \texttt{hug}\;\texttt{s}
    &&\text{2 tokens},\\
  \texttt{bug}  &\rightarrow \texttt{b}\;\texttt{ug}
    &&\text{2 tokens --- never seen in training, still representable},\\
  \texttt{mug}  &\rightarrow \texttt{?}\;\texttt{ug}
    &&\text{\texttt{m} is not in the vocabulary at all.}
\end{aligned}
\]
That last line is the whole reason modern tokenizers are \emph{byte}-level: by
starting from the 256 possible bytes rather than from observed characters, no
input can ever be unrepresentable. The cost is that unusual text fragments into
many short tokens.
\end{workedexample}

\begin{workedexample}[What a context window actually holds]
English averages roughly $0.75$ words per token, so a $4{,}096$-token window
holds about
\[
  4096 \times 0.75 \approx 3{,}070 \text{ words } \approx 12 \text{ pages.}
\]
Conversely, a $10{,}000$-word document needs about
$10{,}000 / 0.75 \approx 13{,}300$ tokens --- it does not fit, and must be
chunked.

\emph{The ratio is not constant.} Code, names, numbers, and non-English text
tokenize far less efficiently. The same sentence in a language written without
spaces, or in one under-represented in the training corpus, can cost two to
three times more tokens than its English equivalent. Since both cost and
context limits are billed in tokens, this is a real and frequently overlooked
fairness issue as well as a budgeting one.
\end{workedexample}

\begin{cautionbox}[The tokenizer explains most ``stupid'' model failures]
When a model cannot count the letters in a word, cannot reverse a string, or
does arithmetic badly on long numbers, the cause is usually the tokenizer, not
the reasoning. The model never saw the letters --- it saw two or three opaque
subword IDs. A number like \texttt{1234567} may split as
\texttt{123}$\vert$\texttt{45}$\vert$\texttt{67}, with no digit-place structure
preserved at all. Before concluding that a model ``cannot reason'', check what
it was actually given.
\end{cautionbox}
\begin{v3topic}{Why Subword Tokenization Dominates}
    Subword tokenization solves three fundamental NLP problems\textsuperscript{\cite[][pp.~15--16]{farrisHowLargeLanguage2025}}:
    \begin{itemize}
        \item \textbf{OOV handling}: unknown words decompose into known subwords rather than failing entirely
        \item \textbf{Morphological awareness}: prefixes and suffixes become reusable tokens (\emph{un-}, \emph{-ness}, \emph{-ing})
        \item \textbf{Multilingual coverage}: any Unicode text can be represented; no whitespace required (works for Chinese, Japanese)
    \end{itemize}
    \textbf{Vocabulary size trade-off}: a larger vocabulary often yields fewer tokens per text and more specific fragments, but enlarges the embedding and output matrices; a smaller vocabulary usually creates longer sequences from more reusable fragments.
    BPE is greedy: ``I'm running'' tokenises to 3 tokens; ``I'm runnin'' (missing \emph{g}) tokenises to 4 tokens because \emph{runnin} is not in the learned vocabulary.
\end{v3topic}

% ---- 12.1.4 Evaluation of NLP Systems ------------------------------------------------------
\subsection{Evaluation of NLP Systems}

\begin{v3topic}{Automatic NLP Metrics}
    NLP systems are evaluated with a mix of automatic metrics and human judgment\textsuperscript{\cite[][pp.~11--12]{dhamaniIntroductionGenerativeAI2024}}:

    \begin{tblr}{width=\linewidth,
        colspec = {Q[3cm]X[2.5]X[4]},
        hlines, vlines,
        row{1} = {font=\bfseries\small, bg=LimeGreen!30},
        rows = {font=\small},
    }
        Metric & Task & Key formula / idea \\
        \textbf{Perplexity} & Language modelling & $\mathrm{PP} = \exp\!\left(-\frac{1}{N}\sum_{i=1}^{N}\log P(w_i|\text{context})\right)$; lower = better \\
        \textbf{BLEU}\textsuperscript{\cite{papineniBLEUMethodAutomatic2002}} & Machine translation & Precision of $n$-gram overlaps between hypothesis and reference; brevity penalty \\
        \textbf{ROUGE-N}\textsuperscript{\cite{linROUGEPackageAutomatic2004}} & Summarisation & Recall of $n$-gram overlaps: $\frac{\text{matched }n\text{-grams}}{\text{reference }n\text{-grams}}$ \\
        \textbf{F1} & NER, QA, classification & Harmonic mean of precision and recall: $\frac{2\,P\,R}{P+R}$ \\
        \textbf{GLUE / SuperGLUE} & General NLU & Multi-task benchmark suite; BERT$_\mathrm{LARGE}$ scores 82.1 \\
        \textbf{Human evaluation} & Fluency, coherence, factuality & Gold standard; expensive; high inter-annotator variance \\
    \end{tblr}
\captionof{table}{The standard automatic metrics. Every one of them measures
surface overlap or model confidence; none of them measures whether the output is
true.}
\label{tab:nlp-metrics}

    Automatic metrics correlate imperfectly with human judgment: a translation can score high on BLEU while being pragmatically awkward.
    Human evaluation remains the gold standard for high-stakes NLP applications.

\begin{workedexample}[Perplexity, computed and interpreted]
A language model assigns these probabilities to the four words of a sentence:
$0.5,\ 0.25,\ 0.25,\ 0.5$.

\emph{Cross-entropy first} (in bits):
\[
  H = -\tfrac{1}{4}\bigl(\log_2 0.5 + \log_2 0.25 + \log_2 0.25
      + \log_2 0.5\bigr)
    = -\tfrac{1}{4}(-1-2-2-1) = \tfrac{6}{4} = 1.5 \text{ bits.}
\]

\emph{Perplexity} is just the exponential of that:
\[
  \mathrm{PP} = 2^{1.5} = 2.83 .
\]

\emph{Read it.} At each word, the model is as uncertain as if it were choosing
uniformly among $2.83$ equally likely options. That is the only correct
intuition for perplexity: it is an \emph{effective branching factor}.

\emph{Two warnings.} Perplexity depends on the tokenizer, so two models with
different vocabularies cannot be compared by perplexity at all. And a model can
have excellent perplexity while being confidently, fluently wrong --- perplexity
measures how well the model predicts text, never whether the text is true.
\end{workedexample}

\begin{workedexample}[A BLEU score, all the way through]
\textsf{Reference:} ``the cat is on the mat'' ($r = 6$ words). \\
\textsf{Candidate:} ``the cat is on mat'' ($c = 5$ words).

\emph{Step 1 --- clipped $n$-gram precisions.}
\[
\begin{aligned}
  p_1 &= 5/5 = 1.000 &&\text{(all five unigrams appear in the reference)},\\
  p_2 &= 3/4 = 0.750 &&\text{(\textsf{on mat} is not a reference bigram)},\\
  p_3 &= 2/3 = 0.667, \qquad p_4 &&= 1/2 = 0.500 .
\end{aligned}
\]

\emph{Step 2 --- geometric mean.}
\[
  \bigl(1.000 \times 0.750 \times 0.667 \times 0.500\bigr)^{1/4}
  = (0.250)^{1/4} = 0.707 .
\]

\emph{Step 3 --- brevity penalty.} The candidate is shorter than the reference
($c < r$), so
\[
  \mathrm{BP} = e^{\,1 - r/c} = e^{\,1-6/5} = e^{-0.2} = 0.819 .
\]

\emph{Step 4 --- BLEU.} $0.819 \times 0.707 = \mathbf{0.58}$.

\emph{What it missed.} Dropping one \textsf{the} cost a third of the score, and
BLEU cannot tell you that the sentence still means exactly the same thing. Now
consider the reverse: ``the cat is not on the mat'' shares five of six words
with the reference and would score very well, while asserting the opposite. BLEU
counts overlapping strings. It has no access to meaning, and no access to truth.
\end{workedexample}

\begin{cautionbox}[Report the metric with the thing it cannot see]
Every automatic metric in \cref{tab:nlp-metrics} is a proxy. BLEU and ROUGE
reward surface overlap and punish valid paraphrase; perplexity rewards
predictability and is blind to factuality; a benchmark score rewards whatever
the benchmark happened to sample, and is vulnerable to that benchmark having
leaked into pre-training data. State the metric, its known blind spot, and the
human or task-level check you ran alongside it. A number without its blind spot
is a marketing claim.
\end{cautionbox}
\end{v3topic}

% ==== Section 2: Text Processing and Representation ==========================================
\section{Text Processing and Representation}

% ---- 12.2.1 Word Vectors and Embeddings ----------------------------------------------------
\subsection{Word Vectors and Embeddings}

\begin{v3topic}{Classical Text Representations}
    Before neural embeddings, text was represented with sparse, count-based vectors\textsuperscript{\cite[][p.~34]{farrisHowLargeLanguage2025}}:
    \begin{itemize}
        \item \textbf{Bag-of-Words (BoW)}: document as a vector of word counts; ignores word order; vocabulary size can exceed 100k dimensions
        \item \textbf{TF-IDF} (Term Frequency -- Inverse Document Frequency): weights terms by how distinctive they are across the corpus:
    \end{itemize}
    \begin{equation}
        \mathrm{tfidf}(t, d) = \underbrace{\frac{\text{count}(t,d)}{|d|}}_{\text{TF}} \times \underbrace{\log\frac{N}{\text{df}(t)}}_{\text{IDF}}
    \end{equation}
\where{$t$ is a term, $d$ a document, $|d|$ its length in tokens, $N$ the
number of documents in the corpus, and $\text{df}(t)$ the number of documents
containing $t$. The logarithm is what suppresses words that appear
everywhere.}
    Common words like ``the'' get low IDF weight; rare but discriminative words get high weight.
    TF-IDF remains useful for information retrieval and keyword extraction but cannot capture semantic similarity.

\end{v3topic}
\begin{figure}[htbp]
\centering
\begin{tikzpicture}[scale=1.0,
  ax/.style={->,draw=AlmanachInk!45,>={Stealth[length=1.6mm]}},
  w/.style={circle,fill=AlmanachTeal,inner sep=1.3pt},
  wo/.style={circle,fill=AlmanachOchre,inner sep=1.3pt},
  lbl/.style={font=\sffamily\scriptsize,color=AlmanachInk!80},
  rel/.style={-{Stealth[length=1.8mm]},draw=AlmanachRed,thick,dashed}]
  \draw[ax] (0,0) -- (7.2,0);
  \draw[ax] (0,0) -- (0,3.4);
  \node[lbl,rotate=90,anchor=south] at (-4mm,1.7) {latent dimension 2};
  \node[lbl,anchor=north] at (3.6,-3mm) {latent dimension 1};

  \node[w] (king) at (1.4,2.5) {};   \node[lbl,above] at (1.4,2.6) {king};
  \node[w] (queen) at (3.4,2.5) {};  \node[lbl,above] at (3.4,2.6) {queen};
  \node[w] (man) at (1.4,1.0) {};    \node[lbl,below] at (1.4,0.85) {man};
  \node[w] (woman) at (3.4,1.0) {};  \node[lbl,below] at (3.4,0.85) {woman};
  \draw[rel] (man) -- (woman);
  \draw[rel] (king) -- (queen);
  \node[lbl,color=AlmanachRed] at (2.4,1.22) {same offset};

  \node[wo] (cat) at (5.7,2.2) {};   \node[lbl,right=1mm] at (5.75,2.2) {cat};
  \node[wo] (dog) at (6.1,2.6) {};   \node[lbl,right=1mm] at (6.15,2.6) {dog};
  \node[wo] (kitten) at (5.4,2.6) {};\node[lbl,left=1mm] at (5.35,2.6) {kitten};
  \draw[draw=AlmanachOchre,dashed] (5.75,2.45) circle (0.62);
  \node[lbl,color=AlmanachOchre,anchor=north] at (5.75,1.72)
    {similar words cluster};
\end{tikzpicture}
\caption{What an embedding buys. Words used in similar contexts end up near
each other, and some relations appear as consistent offsets --- the basis of
the famous \emph{king $-$ man $+$ woman $\approx$ queen} arithmetic. The
regularity is real but weaker and more selective than the popular examples
suggest, and it holds far better for frequent words than for rare ones.}
\label{fig:nlp-embedding-space}
\end{figure}

\begin{v3topic}{Why the Dimensions Have to Be Many}
Two dimensions suffice to separate cats from cars. But a single token must
simultaneously carry that \emph{bank} may be a financial institution or a river
edge, that \emph{run} may be a verb or a noun, and that \emph{Java} may be a
language or an island. Each dimension can encode a different aspect of usage,
and a token's position encodes all of its potential meanings at once ---
context later selects among them.

\begin{engineernote}
An embedding table is a $V\times d$ array and a lookup is a single memory read
of $d$ numbers --- no arithmetic at all. It is also large: at $V = 50{,}000$ and
$d = 4{,}096$ in 16-bit precision the table alone is
$50{,}000 \times 4{,}096 \times 2 = 400$\,MB. That makes it the one component
of a model that is already cheap in operations and expensive in memory, which
is exactly the wrong way round for the memory wall described in
\cref{ch:deep-learning}.
\end{engineernote}

\begin{cautionbox}[Embeddings are inherently continuous]
The difference between \emph{happy} and \emph{joyful} is a small displacement
in this space. Rounding each dimension to a single bit would destroy precisely
the fine distinctions the representation exists to carry --- which is why
aggressive quantisation is usually applied to the weight matrices rather than
to the embeddings, and why embedding tables are often kept at higher precision
than the rest of the model.
\end{cautionbox}
\end{v3topic}

\begin{v3topic}{How Meaning Survives Subword Splitting}
If embeddings carry meaning and the vocabulary is made of fragments, an obvious
objection follows: BPE may split \emph{prince} into \texttt{prin} + \texttt{ce},
and neither piece means anything.

The resolution is that the embedding table only supplies raw material; the
Transformer layers compose meaning from it.

\begin{enumerate}
  \item \textbf{Lookup.} \texttt{prin} and \texttt{ce} each receive their own
    vector. After training, \texttt{prin} encodes whatever
    \emph{prince}, \emph{print}, and \emph{principal} have in common --- a
    partial, ambiguous signal.
  \item \textbf{Attention assembles the word.} In the first layer, \texttt{prin}
    attends to \texttt{ce} and vice versa, and each representation absorbs the
    other.
  \item \textbf{Depth completes it.} After a few layers the vector that began
    as \texttt{prin} effectively represents \emph{prince in this sentence}. By
    the deeper layers the representations are not about subword pieces at all.
\end{enumerate}

\begin{plainlanguage}
The model does not distinguish between assembling a word out of fragments and
assembling a sentence out of words. Both are the same operation applied at a
different scale, which is why an unfamiliar word, a misspelling, or a line of
source code is not a special case for a language model --- \texttt{princ} +
\texttt{ess} composes exactly as \texttt{prin} + \texttt{ce} does.
\end{plainlanguage}

\begin{engineernote}
This also explains a limit. Because composition happens \emph{in the layers},
a fragment's meaning depends on there being enough layers above it to do the
assembling. It is one reason why heavily fragmented text --- an
under-represented language, an unusual identifier, a long number --- is handled
less well: the model spends its early layers reconstructing what a
better-matched tokenizer would have handed it for free.
\end{engineernote}
\end{v3topic}

\begin{v3topic}{Word2Vec: Learning Dense Embeddings}
    Word2Vec is due to Mikolov and
    colleagues\textsuperscript{\cite{mikolovEfficientEstimationWord2013}}.

    \textbf{Word2Vec} (Mikolov et al., Google, 2013) learns dense, low-dimensional word vectors from large corpora using two architectures\textsuperscript{\cite[][p.~36]{farrisHowLargeLanguage2025}}:
    \begin{itemize}
        \item \textbf{CBOW} (Continuous Bag-of-Words): predict center word from surrounding context words
        \item \textbf{Skip-gram}: predict surrounding context words given the center word
    \end{itemize}
    Training signal: maximise $P(\text{context}|\text{word})$ via negative sampling.
    Result: 300-dimensional vectors that encode semantic relationships geometrically.
    The famous analogy holds:
    \begin{equation}
        \vec{v}(\text{king}) - \vec{v}(\text{man}) + \vec{v}(\text{woman}) \approx \vec{v}(\text{queen})
    \end{equation}
\where{$\vec{v}(w)$ is the embedding vector of word $w$; the relation holds
approximately and far better for frequent words than rare ones.}
    Word2Vec vectors are \emph{static}: the same vector is assigned to ``bank'' regardless of context.

\end{v3topic}
\begin{v3topic}{GloVe and the Embedding Geometry}
    GloVe factorises a global co-occurrence
    matrix\textsuperscript{\cite{penningtonGloVeGlobalVectors2014}}.

    \textbf{GloVe} (Global Vectors, Pennington et al., Stanford 2014) learns embeddings from global co-occurrence statistics rather than local context windows\textsuperscript{\cite[][p.~35]{farrisHowLargeLanguage2025}}.
    Objective: $\vec{v}_i \cdot \vec{v}_j + b_i + b_j = \log X_{ij}$ where $X_{ij}$ is the co-occurrence count.
    GloVe captures corpus-wide statistics more efficiently than Word2Vec on large corpora.
\visualseparator
    \textbf{Embedding space interpretation}: synonyms cluster nearby; analogies follow vector arithmetic.
    Words with similar meaning are close in cosine similarity.
    \emph{Limitation}: static embeddings assign a single vector per word token, ignoring polysemy --- ``bank'' (finance) and ``bank'' (river) get the same representation.

\end{v3topic}
\begin{v3topic}{Contextual Embeddings}
    Static embeddings (Word2Vec, GloVe) fail for polysemous words.
    \textbf{Contextual embeddings} produce a different vector for the same word depending on its sentence context\textsuperscript{\cite[][p.~35]{farrisHowLargeLanguage2025}}:
    \begin{itemize}
        \item \textbf{ELMo} (Peters et al., 2018): bidirectional LSTM; concatenates forward and backward representations
        \item \textbf{BERT} embeddings: transformer-based; full bidirectional context at every layer
        \item \textbf{GPT embeddings}: left-to-right context only (causal)
    \end{itemize}
    Modern NLP systems no longer treat embeddings as a pre-processing step; the entire transformer is fine-tuned end-to-end, and the contextual representations emerge from pre-training on massive corpora.
\end{v3topic}

% ---- 12.2.2 Statistical and Classical Approaches ------------------------------------------
\subsection{Statistical and Classical Approaches}

\begin{v3topic}{N-gram Language Models}
    An \textbf{$N$-gram language model} estimates word probabilities using the Markov assumption --- only the previous $N-1$ words matter\textsuperscript{\cite[][]{bengioNeuralProbabilisticLanguage}}:
    \begin{equation}
        P(w_t \mid w_1, \ldots, w_{t-1})
        \approx
        P(w_t \mid w_{t-N+1}, \ldots, w_{t-1})
    \end{equation}
\where{$w_t$ is the word at position $t$ and $N$ the order of the model, so
only the previous $N-1$ words are retained. Larger $N$ captures more context and
makes unseen $n$-grams exponentially more likely.}
    Counts in a training corpus estimate these probabilities.
    Smoothing (Laplace, Kneser-Ney) handles unseen $n$-grams.
    Bengio et al.\ (2003) introduced the \textbf{neural probabilistic language model}: replace the count-based estimate with a feed-forward neural network that maps word embeddings to the next-word distribution.
    This was the conceptual precursor to modern LLMs.

\begin{workedexample}[Counting bigrams, and the zero that breaks everything]
Corpus: \textsf{``the cat sat on the mat the cat ran''} --- nine tokens, six
distinct words $V = \{$the, cat, sat, on, mat, ran$\}$.

\emph{Raw counts.} \textsf{the} occurs $3$ times. It is followed by
\textsf{cat} twice and \textsf{mat} once. So
\[
  P(\textsf{cat}\mid\textsf{the}) = \tfrac{2}{3} = 0.667,
  \qquad
  P(\textsf{mat}\mid\textsf{the}) = \tfrac{1}{3} = 0.333,
  \qquad
  P(\textsf{sat}\mid\textsf{the}) = \tfrac{0}{3} = \mathbf{0}.
\]

That zero is fatal. Sentence probability is a product, so a single unseen
bigram makes the \emph{entire sentence} impossible --- and its log-probability
$-\infty$.

\emph{Add-one (Laplace) smoothing.} Add $1$ to every count and $|V| = 6$ to
every denominator:
\[
  P(\textsf{cat}\mid\textsf{the}) = \tfrac{2+1}{3+6} = \tfrac{3}{9} = 0.333,
  \qquad
  P(\textsf{sat}\mid\textsf{the}) = \tfrac{0+1}{3+6} = \tfrac{1}{9} = 0.111 .
\]
Check that it still sums to one over $V$:
$\tfrac{1+3+1+1+2+1}{9} = \tfrac{9}{9} = 1$. \checkmark

\emph{The cost.} The observed bigram fell from $0.667$ to $0.333$ --- add-one
smoothing took half the probability mass away from what we actually saw and
handed it to things we never saw. With a realistic vocabulary of $50{,}000$
words the denominator becomes $3 + 50{,}000$ and the observed evidence is
annihilated entirely. This is why real systems use Kneser--Ney or backoff rather
than Laplace, and ultimately why neural models --- which share statistical
strength through embeddings instead of counting --- replaced $n$-grams
altogether.
\end{workedexample}

\end{v3topic}
\begin{v3topic}{Sequence Models: HMM and CRF}
    Classical NLP used probabilistic sequence models for structured prediction (POS tagging, NER)\textsuperscript{\cite[][p.~6]{dhamaniIntroductionGenerativeAI2024}}:
    \begin{itemize}
        \item \textbf{Hidden Markov Model (HMM)}: generative model; hidden states = POS tags; observed outputs = words.
              Computes $P(\text{tags}|\text{words})$ via Viterbi decoding.
        \item \textbf{Naive Bayes}: fast generative classifier that models class priors and class-conditional features;
              $P(c \mid \mathbf{x}) \propto P(c) \prod_i P(x_i \mid c)$.
              Still competitive for sentiment analysis with sparse features.
        \item \textbf{Conditional Random Field (CRF)}: discriminative sequence model; models $P(\mathbf{y}|\mathbf{x})$ directly; avoids HMM's independence assumptions.
              State-of-the-art for NER before transformers.
    \end{itemize}
    These models are now largely superseded by transformer-based approaches but remain important baselines and are still used in low-resource settings.
\end{v3topic}

% ---- 12.2.3 RNN-based Approaches ----------------------------------------------------------
\subsection{RNN-based Approaches}

\begin{v3topic}{Seq2Seq: Encoder-Decoder for NLP}
    The encoder--decoder formulation is due to Sutskever, Vinyals and
    Le\textsuperscript{\cite{sutskeverSequenceSequenceLearning2014}}.

    The \textbf{sequence-to-sequence (seq2seq)} architecture (Sutskever et al., 2014) enables NLP tasks with variable-length inputs and outputs\textsuperscript{\cite[][p.~8]{dhamaniIntroductionGenerativeAI2024}}:
    \begin{enumerate}
        \item \textbf{Encoder} (LSTM/GRU): reads input sequence and compresses it into a fixed context vector $\mathbf{c}$
        \item \textbf{Decoder} (LSTM/GRU): autoregressively generates output conditioned on $\mathbf{c}$ and previously generated tokens
    \end{enumerate}
    Applications: machine translation, summarisation, dialogue, question answering.
    \textbf{Limitation}: the fixed-size context vector becomes a bottleneck for long sequences; the encoder must compress arbitrary input length into a single vector, losing detail.

\end{v3topic}
\begin{v3topic}{Attention in RNNs}
    Attention was introduced for translation by Bahdanau, Cho and
    Bengio\textsuperscript{\cite{bahdanauNeuralMachineTranslation2015}}, four
    years before it displaced recurrence entirely
    (\cref{sec:dl-attention}).

    The \textbf{attention mechanism} (Bahdanau et al., 2014) overcomes the seq2seq bottleneck by allowing the decoder to selectively attend to all encoder hidden states at each generation step\textsuperscript{\cite[][p.~8]{dhamaniIntroductionGenerativeAI2024}}:
    \begin{equation}
        \mathbf{c}_t = \sum_{i} \alpha_{t,i}\, \mathbf{h}_i, \quad \alpha_{t,i} = \mathrm{softmax}(e_{t,i})
    \end{equation}
    where $e_{t,i} = \mathrm{score}(\mathbf{s}_{t-1}, \mathbf{h}_i)$ is an alignment score between the current decoder state $\mathbf{s}_{t-1}$ and encoder hidden state $\mathbf{h}_i$.
    The model learns \emph{which} source tokens to focus on when generating each target token.
    This was a major breakthrough for machine translation and inspired the transformer's self-attention.
    \textbf{Key limit}: RNNs are inherently sequential and cannot be parallelised --- making them slow to train on large corpora.
\end{v3topic}

% ---- 12.2.4 Transformer-based Approaches ---------------------------------------------------
\subsection{Transformer-based Approaches}

\begin{v3topic}{The Transformer for NLP}
    The \textbf{Transformer} (Vaswani et al., 2017) replaced recurrence with \emph{self-attention}, enabling full parallelisation during training\textsuperscript{\cite[][]{vaswaniAttentionAllYou2023}}.
    See Ch.~6 for the full architecture. The self-attention operator is:
    \begin{equation}
        \mathrm{att}(Q, K, V) = \mathrm{softmax}\!\left(\frac{QK^\top}{\sqrt{D^{\mathrm{qk}}}}\right)\!V
    \end{equation}
    where $Q$, $K$, $V$ are queries, keys, values derived from the input\textsuperscript{\cite[][eq.~4.3]{fleuretLittleBookDeep}}.
    Three transformer variants serve different NLP tasks:
    \begin{itemize}
        \item \textbf{Encoder-only} (BERT, RoBERTa): bidirectional context; best for classification, NER, QA
        \item \textbf{Decoder-only} (GPT, Gemini): causal/autoregressive; best for generation
        \item \textbf{Encoder-decoder} (T5, Google Translate): seq2seq tasks (translation, summarisation)
    \end{itemize}

\end{v3topic}
\begin{v3topic}{BERT: Pre-training for NLU}
    \textbf{BERT} (Bidirectional Encoder Representations from Transformers; Devlin et al., 2019) introduced self-supervised pre-training for language understanding\textsuperscript{\cite[][pp.~1--4]{devlinBERTPretrainingDeep2019}}:

    \textbf{Pre-training tasks} (unsupervised, on BooksCorpus + Wikipedia):
    \begin{itemize}
        \item \textbf{Masked Language Model (MLM)}: randomly mask 15\% of input tokens; predict masked tokens from bidirectional context
        \item \textbf{Next Sentence Prediction (NSP)}: classify whether sentence B follows sentence A
    \end{itemize}
    \textbf{Tokenizer}: WordPiece; vocabulary size 30{,}000.
    \textbf{Scale}: BERT$_\mathrm{BASE}$ (110M parameters, 12 layers); BERT$_\mathrm{LARGE}$ (340M parameters, 24 layers).
    \textbf{Fine-tuning}: add a single task-specific output layer; fine-tune all parameters end-to-end.
    Result: GLUE score 82.1 (BERT$_\mathrm{LARGE}$) vs 74.0 (pre-BERT SOTA) --- a landmark result\textsuperscript{\cite[][p.~5]{devlinBERTPretrainingDeep2019}}.

\end{v3topic}
\begin{v3topic}{GPT: Pre-training for Generation}
    The \textbf{GPT} family (Generative Pre-trained Transformers; OpenAI, 2018--) uses decoder-only transformers trained autoregressively to predict the next token\textsuperscript{\cite[][pp.~8--9]{dhamaniIntroductionGenerativeAI2024}}:
    \begin{itemize}
        \item \textbf{GPT-1} (2018): pre-trained on BooksCorpus; fine-tuned per task
        \item \textbf{GPT-2} (2019): 1.5B parameters; delayed release due to misuse concerns
        \item \textbf{GPT-3} (2020): 175B parameters; \emph{few-shot in-context learning} without fine-tuning
        \item \textbf{GPT-4} (2023): multimodal; state-of-art on most NLP benchmarks\textsuperscript{\cite[][]{openaiGPT4TechnicalReport2023}}
    \end{itemize}
    The key insight: pre-training on massive unlabelled text followed by fine-tuning or prompting is far more data-efficient than training task-specific models from scratch.
    BERT (encoder) and GPT (decoder) are complementary: BERT excels at understanding; GPT at generation.
\end{v3topic}

% ==== Section 3: Speech Processing ===========================================================
\section{Speech Processing}

% ---- 12.3.1 Speech Recognition (ASR) -------------------------------------------------------
\subsection{Speech Recognition (ASR)}

\begin{v3topic}{Classical ASR Pipeline}
    Classical \textbf{Automatic Speech Recognition (ASR)} decomposes recognition into three components combined via Bayes' theorem\textsuperscript{\cite[][p.~11]{dhamaniIntroductionGenerativeAI2024}}:
    \begin{equation}
        P(\mathbf{w} \mid \mathbf{a}) \propto \underbrace{P(\mathbf{a} \mid \mathbf{w})}_{\text{acoustic model}} \times \underbrace{P(\mathbf{w})}_{\text{language model}}
    \end{equation}
\where{$\mathbf{a}$ is the observed acoustic signal and $\mathbf{w}$ a
candidate word sequence; the acoustic model scores how well the sounds match the
words, the language model how plausible the word sequence is on its own.}
    \begin{itemize}
        \item \textbf{Acoustic model}: maps audio features (MFCC) to phoneme probabilities; classically HMM-GMM, later deep neural networks
        \item \textbf{Language model}: $n$-gram or neural LM providing prior over word sequences
        \item \textbf{Pronunciation lexicon}: maps phoneme sequences to words
    \end{itemize}
    These components were trained separately and combined at inference.
    \textbf{CTC (Connectionist Temporal Classification)}: end-to-end objective (Graves et al., 2006) that allows the network to produce output sequences without frame-by-frame alignment; introduces a blank symbol to handle alignment ambiguity.

\end{v3topic}
\begin{v3topic}{End-to-End ASR: Whisper}
    Whisper is trained with large-scale weak
    supervision\textsuperscript{\cite{radfordRobustSpeechRecognition2023}}.

    \textbf{Whisper} (Radford et al., OpenAI 2022) is a transformer encoder-decoder trained end-to-end on 680{,}000 hours of diverse, weakly labelled multilingual audio\textsuperscript{\cite[][p.~11]{dhamaniIntroductionGenerativeAI2024}}:
    \begin{itemize}
        \item \textbf{Encoder}: processes 30-second log-mel spectrogram windows; generates contextual audio representations
        \item \textbf{Decoder}: autoregressively generates text tokens conditioned on encoder output
        \item No separate acoustic model, language model, or lexicon required
        \item Handles 99 languages, punctuation, timestamps, and language identification natively
    \end{itemize}
    Whisper demonstrates that scale and diversity of training data can replace carefully engineered pipeline components.
    It is robust to accents, background noise, and domain shift to a degree that classical pipelines struggle to achieve.
\end{v3topic}

% ---- 12.3.2 Speech Synthesis (TTS) ---------------------------------------------------------
\subsection{Speech Synthesis (TTS)}

\begin{v3topic}{Classical and Neural TTS}
    \textbf{Text-to-Speech (TTS)} converts written text into natural-sounding audio.
    Three generations of TTS systems exist\textsuperscript{\cite[][p.~11]{dhamaniIntroductionGenerativeAI2024}}:
    \begin{itemize}
        \item \textbf{Concatenative TTS}: stitch together pre-recorded speech units (diphones, triphones) from a large database; high naturalness but inflexible; requires new recordings per voice
        \item \textbf{Parametric TTS (HMM-based)}: generate speech by predicting vocal tract parameters from text via HMM; flexible but robotic quality
        \item \textbf{Neural TTS}: end-to-end neural synthesis; dramatically improves naturalness
    \end{itemize}

\end{v3topic}
\begin{v3topic}{Neural TTS: WaveNet and Tacotron}
    WaveNet models raw audio
    autoregressively\textsuperscript{\cite{oordWaveNetGenerativeModel2016}}.

    Neural TTS systems operate in two stages: text $\to$ acoustic features $\to$ waveform\textsuperscript{\cite[][p.~11]{dhamaniIntroductionGenerativeAI2024}}:
    \begin{itemize}
        \item \textbf{WaveNet} (van den Oord et al., DeepMind, 2016): dilated causal convolutional network that generates raw audio waveforms sample-by-sample at 16kHz; produces natural prosody and voice quality; very slow at inference (autoregressive)
        \item \textbf{Tacotron 2} (Google, 2018): seq2seq RNN that maps text to mel spectrograms; coupled with a WaveNet vocoder to synthesise audio from the spectrogram
        \item \textbf{FastSpeech 2}: transformer-based non-autoregressive model; fast inference, controllable speaking rate and pitch
        \item \textbf{Modern vocoders}: diffusion-based (DiffWave) or GAN-based (HiFi-GAN) waveform synthesisers replace WaveNet for real-time applications
    \end{itemize}
\end{v3topic}

% ==== Section 4: Large Language Models and Generative AI =====================================
\section{Large Language Models and Generative AI}

% ---- 12.4.1 Pre-training and Fine-tuning ----------------------------------------------------
\subsection{Autoregressive Generation: One Token at a Time}
\label{sec:nlp-autoregressive-generation}

\Cref{ch:deep-learning} built the Transformer, and the sections above turned
text into tokens and tokens into vectors. One step is still missing, and it is
the one that turns all of that into a system that \emph{writes}: how a stack of
matrices becomes a stream of words.

\begin{v3topic}{From the Final Vector to the Next Token}
After the last Transformer block, each position holds a $d$-dimensional vector.
The \keyterm{output head} --- a single matrix of shape $d \times V$ --- projects
the vector at the \emph{final} position into one score per vocabulary entry:
\begin{equation}
  \mathbf{z} = W_{\text{vocab}}\,\mathbf{h}_T \in \mathbb{R}^{V},
  \qquad
  p = \mathrm{softmax}(\mathbf{z}) .
  \label{eq:nlp-output-head}
\end{equation}
\where{$\mathbf{h}_T$ is the final-layer vector at the last position, $V$ the
vocabulary size, and $\mathbf{z}$ the vector of \keyterm{logits} --- one raw
score per possible next token. Many models \emph{tie} $W_{\text{vocab}}$ to the
embedding table, reusing the same weights and paying for them once.}

The loop is then: pick a token from $p$, append it to the input, and run the
whole model again. That is the entire generation algorithm.

\begin{figure}[htbp]
\centering
\begin{tikzpicture}[
  node distance=5mm,
  bx/.style={draw=AlmanachTeal,fill=AlmanachMist,rounded corners=1mm,
    minimum height=8mm,minimum width=20mm,font=\sffamily\scriptsize,
    align=center},
  pick/.style={bx,fill=AlmanachOchre!22,draw=AlmanachOchre},
  arr/.style={-{Stealth[length=1.7mm]},thick,draw=AlmanachInk!55},
  lbl/.style={font=\sffamily\scriptsize,color=AlmanachInk!70}]
  \node[bx] (ctx) {tokens so far};
  \node[bx,right=of ctx] (model) {Transformer\\stack};
  \node[bx,right=of model] (head) {output head\\\scriptsize $d\to V$};
  \node[bx,right=of head] (soft) {softmax};
  \node[pick,right=of soft] (samp) {sample};
  \draw[arr] (ctx) -- (model); \draw[arr] (model) -- (head);
  \draw[arr] (head) -- (soft); \draw[arr] (soft) -- (samp);
  \draw[arr,draw=AlmanachOchre] (samp.south) -- ++(0,-7mm) -| (ctx.south)
    node[near start,below,lbl,color=AlmanachOchre]
    {append the chosen token, then run the whole model again};
\end{tikzpicture}
\caption{The generation loop. Every token costs one full forward pass through
every layer, which is why generation is irreducibly sequential and why the
key--value cache of \cref{ch:deep-learning} exists --- it avoids recomputing the
history, but it cannot remove the loop.}
\label{fig:nlp-generation-loop}
\end{figure}

\begin{cautionbox}[Why the model may not look ahead]
During training the model sees whole sequences at once, which would let
position $t$ attend to position $t+1$ and read the answer it is being asked to
predict. The \keyterm{causal mask} prevents this by setting every score above
the diagonal to $-\infty$ \emph{before} the softmax, so those weights become
exactly zero (\cref{fig:dl-attention-matrix}).

Two consequences worth holding on to. Masking before rather than after the
softmax is what makes the remaining weights sum to one over the permitted
positions only --- masking afterwards would leave them normalised over
positions that were then discarded. And the mask is why a decoder can be
trained on all positions in parallel while still only ever being generated one
token at a time: training is parallel, inference is not.
\end{cautionbox}
\end{v3topic}

\begin{v3topic}{Choosing the Next Token}
The distribution $p$ does not decide anything on its own --- a \keyterm{decoding
strategy} does, and it is a genuine engineering choice with a large effect on
output.

\begin{workedexample}[Temperature, top-$k$, and top-$p$ on the same logits]
Five candidate tokens with logits $\mathbf{z} = (3.0,\ 2.0,\ 1.0,\ 0.5,\ 0.0)$.

\emph{Plain softmax} ($T=1$): exponentials $20.09,\ 7.39,\ 2.72,\ 1.65,\ 1.00$
sum to $32.84$, giving
\[
  p = (0.612,\ 0.225,\ 0.083,\ 0.050,\ 0.031).
\]

\emph{Temperature} divides the logits by $T$ before the softmax:
\[
\begin{aligned}
  T = 0.5:\quad & p = (0.860,\ 0.116,\ 0.016,\ 0.006,\ 0.002)
    &&\text{sharper --- nearly deterministic},\\
  T = 2.0:\quad & p = (0.403,\ 0.244,\ 0.148,\ 0.115,\ 0.090)
    &&\text{flatter --- more adventurous}.
\end{aligned}
\]
$T\to0$ approaches always picking the top token; $T\to\infty$ approaches
uniform noise.

\emph{Top-$k$} keeps the $k$ highest and renormalises. With $k=2$ at $T=1$, the
surviving mass is $0.612+0.225 = 0.837$, so
\[
  p' = (0.731,\ 0.269).
\]

\emph{Top-$p$} (nucleus) keeps the smallest set whose cumulative probability
reaches $p$. With $p = 0.9$: the running total is
$0.612 \to 0.837 \to 0.919$, so three tokens are kept and renormalised by
$0.919$:
\[
  p' = (0.665,\ 0.245,\ 0.090).
\]

\emph{The difference that matters.} Top-$k$ always keeps exactly $k$ candidates
--- including when the model is certain, where it forces three bad options to
stay in play, and when the model is genuinely uncertain, where it may cut off
good ones. Top-$p$ adapts: it keeps few candidates when the distribution is
peaked and many when it is flat. That is why nucleus sampling is the usual
default, and why quoting a temperature without the truncation method describes
only half the decoder.
\end{workedexample}

\begin{engineernote}
Set the decoder from the task, not from a default. Extraction, classification,
and code that must compile want greedy decoding or a very low temperature ---
reproducibility is worth more than variety, and greedy decoding is the only
setting that is deterministic. Open-ended writing wants nucleus sampling around
$T \approx 0.7$--$1.0$. And when you report an evaluation result, report the
decoding settings with it: the same model at $T = 0$ and $T = 1$ is, for
measurement purposes, two different systems.
\end{engineernote}
\end{v3topic}

\begin{checkpoint}
A model generates a $500$-token answer from a $2{,}000$-token prompt. State how
many forward passes through the full stack are required, why the key--value
cache does not reduce that number, and what it does reduce instead.
\end{checkpoint}

\subsection{Pre-training and Fine-tuning}


\begin{v3topic}{Pre-training Objectives}
    LLMs are initialised by \textbf{pre-training} on vast unlabelled corpora using
    \emph{self-supervised} objectives that require no manual
    annotation.\textsuperscript{\cite[][p.~9]{dhamaniIntroductionGenerativeAI2024}}

    Two dominant objectives:
    \begin{itemize}
        \item \textbf{Masked Language Modelling (MLM):} randomly mask tokens; predict
              the original from bidirectional context --- used by BERT.
        \item \textbf{Causal (autoregressive) LM:} predict the \emph{next} token given
              all preceding tokens --- used by GPT.
    \end{itemize}

    Both are \emph{pretext tasks}: solving them forces the model to build rich internal
    representations of grammar, facts, and
    reasoning.\textsuperscript{\cite[][p.~5]{raschkaBuildLargeLanguage2025}}

\end{v3topic}
\begin{v3topic}{Transfer Learning}
    \textbf{Transfer learning} reuses representations learnt during pre-training for new
    downstream tasks. The pre-trained model acts as a universal \emph{feature extractor}
    of language.\textsuperscript{\cite[][p.~9]{dhamaniIntroductionGenerativeAI2024}}

    Key insight: a model trained to predict masked or next tokens on billions of words
    implicitly learns syntax, semantics, world knowledge, and reasoning patterns. These
    representations \emph{generalise} to tasks the model was never explicitly trained on.

    Fine-tuning then adapts these representations to a specific task with a small labelled
    dataset, dramatically reducing the amount of task-specific data required compared to
    training from
    scratch.\textsuperscript{\cite[][p.~70]{farrisHowLargeLanguage2025}}

\end{v3topic}
\begin{v3topic}{Pre-training vs.\ Fine-tuning vs.\ Prompting}
    \begin{tblr}{width=\linewidth,
        colspec={X[l,2] X[l,3] X[l,3]},
        hlines, vlines,
        row{1}={font=\bfseries\small, bg=LimeGreen!20},
        row{2-Z}={font=\small},
    }
        Strategy & Description & When to Use \\
        Pre-training & Train on large unlabelled corpus from scratch &
            Foundation model; requires enormous compute \\
        Fine-tuning (SFT) & Update all (or most) weights on labelled task data &
            Sufficient labelled data; dedicated deployment \\
        Few-shot prompting & Provide $k$ examples in context; no weight update &
            Limited labelled data; rapid prototyping \\
        Zero-shot prompting & Task description only; no examples &
            Testing generalisation; minimal setup \\
        RAG & Retrieve relevant documents; inject into prompt &
            Knowledge-intensive tasks; up-to-date information \\
    \end{tblr}
\captionof{table}{Three ways to adapt a model, in decreasing order of cost and increasing order of iteration speed.}
\label{tab:nlp-adaptation-strategies}
    \vspace{2pt}
    \textsuperscript{\cite[][p.~70]{farrisHowLargeLanguage2025}}
    \textsuperscript{\cite[][p.~5]{raschkaBuildLargeLanguage2025}}

\end{v3topic}

% ---- 12.4.2 BERT ---------------------------------------------------------------------------
\subsection{BERT}


\begin{v3topic}{BERT Architecture}
    \textbf{BERT} (\emph{Bidirectional Encoder Representations from Transformers}) is a
    Transformer \emph{encoder} that conditions on \textbf{both left and right context} at
    every layer --- unlike GPT, which is left-to-right
    only.\textsuperscript{\cite[][p.~2]{devlinBERTPretrainingDeep2019}}

    Every input sequence begins with a special \texttt{[CLS]} token; its final hidden
    state $C \in \mathbb{R}^{H}$ serves as the aggregate \emph{sentence representation}
    for classification. A \texttt{[SEP]} token separates paired sentences.

    Input embedding $=$ token emb.~$+$ segment emb.~$+$ position emb.

    \begin{tblr}{width=\linewidth,colspec={X[l]X[l]X[l]}, hlines, vlines,
                 row{1}={font=\bfseries\scriptsize, bg=LimeGreen!20},
                 row{2-Z}={font=\scriptsize}}
        & BERT\textsubscript{BASE} & BERT\textsubscript{LARGE} \\
        Layers $L$ & 12 & 24 \\
        Hidden $H$ & 768 & 1024 \\
        Heads $A$ & 12 & 16 \\
        Parameters & 110M & 340M \\
    \end{tblr}
\captionof{table}{The two BERT configurations. The scaling here is modest by current standards, which is worth remembering when reading the results it produced.}
\label{tab:nlp-bert-sizes}

\end{v3topic}
\begin{v3topic}{Pre-training: MLM and NSP}
    BERT is pre-trained with two unsupervised
    tasks:\textsuperscript{\cite[][p.~3-4]{devlinBERTPretrainingDeep2019}}

    \textbf{Task 1 --- Masked LM (MLM):} 15\,\% of WordPiece tokens are selected at
    random. Of these, 80\,\% are replaced with \texttt{[MASK]}, 10\,\% with a random
    token, and 10\,\% left unchanged. The model predicts the original vocabulary item
    from bidirectional context, forcing it to build a \emph{distributional contextual
    representation} of every token.

    \textbf{Task 2 --- Next Sentence Prediction (NSP):} Given sentence pairs $(A, B)$,
    predict \texttt{IsNext} (50\,\%) or \texttt{NotNext} (50\,\%). This teaches
    inter-sentence relationships, benefiting QA and NLI.

    Pre-training data: BooksCorpus (800M words) + English Wikipedia (2,500M words).

\end{v3topic}
\begin{v3topic}{Fine-tuning BERT for Downstream Tasks}
    Fine-tuning is straightforward: add one task-specific output layer, then fine-tune
    \emph{all} parameters end-to-end on labelled
    data.\textsuperscript{\cite[][p.~4-5]{devlinBERTPretrainingDeep2019}}

    \begin{tblr}{width=\linewidth,
        colspec={X[l,2] X[l,2] X[l,4]},
        hlines, vlines,
        row{1}={font=\bfseries\small, bg=LimeGreen!20},
        row{2-Z}={font=\small},
    }
        Task Type & Input & Output \\
        Sentence classification & \texttt{[CLS]} + sentence & Softmax over classes on $C$ \\
        Sentence-pair classification & \texttt{[CLS]} + sent.~A + \texttt{[SEP]} + sent.~B & Softmax on $C$ \\
        QA (SQuAD) & \texttt{[CLS]} + question + \texttt{[SEP]} + passage & Start/end token positions \\
        Named Entity Recognition & \texttt{[CLS]} + sentence & Per-token label from $T_i$ \\
    \end{tblr}
\captionof{table}{How different task types map onto the same encoder: only the input formatting and the output head change.}
\label{tab:nlp-bert-finetuning}

    BERT achieved state-of-the-art on 11 NLP tasks at release, including GLUE
    (+7.7\,\%), SQuAD~1.1 F1~93.2, and CoNLL-2003 NER F1~92.8.

\end{v3topic}

% ---- 12.4.3 GPT Family ----------------------------------------------------------------------
\subsection{GPT Family}


\begin{v3topic}{GPT: Decoder-Only Architecture}
    GPT (\emph{Generative Pretrained Transformer}) is a \textbf{decoder-only} Transformer
    using \emph{causal} (left-to-right) self-attention: each token can only attend to
    preceding tokens.\textsuperscript{\cite[][p.~12]{raschkaBuildLargeLanguage2025}}

    Generation is \textbf{autoregressive}: the model produces one token at a time,
    appending each output to the input before predicting the next. The pre-training
    objective maximises the conditional log-likelihood:
    \begin{equation}
        \ell(\theta) = \sum_{t=1}^{T} \log P(x_t \mid x_1, \ldots, x_{t-1};\,\theta)
    \end{equation}
\where{$x_t$ is the token at position $t$, $T$ the sequence length, and
$\theta$ the model parameters. This is the entire pre-training objective:
predict each token from everything before it.}
    Equivalently, training minimises the negative log-likelihood $\mathcal{L}(\theta)=-\ell(\theta)$. This is next-token prediction on large unlabelled
    corpora.\textsuperscript{\cite[][p.~129]{raschkaBuildLargeLanguage2025}}

\end{v3topic}
\begin{v3topic}{GPT Evolution and Scaling Laws}
    \begin{tblr}{width=\linewidth,
        colspec={X[l] X[l] X[l]},
        hlines, vlines,
        row{1}={font=\bfseries\scriptsize, bg=LimeGreen!20},
        row{2-Z}={font=\scriptsize},
    }
        Model & Year & Scale / Notes \\
        GPT-1 & 2018 & 117M params; fine-tuned on downstream tasks \\
        GPT-2 & 2019 & 1.5B params; release withheld (misuse risk) \\
        GPT-3 & 2020 & 175B params; few-shot in-context learning \\
        Chinchilla & 2022 & 70B params; compute-optimal parameter--data balance \\
        GPT-4 & 2023 & Multimodal; strong reasoning and coding \\
    \end{tblr}
\captionof{table}{The GPT series by scale. Time-sensitive: treat parameter counts and dates as of publication, and verify before quoting.}
\label{tab:nlp-gpt-evolution}
    \vspace{2pt}
    \textsuperscript{\cite[][p.~9]{dhamaniIntroductionGenerativeAI2024}}

    \textbf{Scaling laws} describe predictable power-law relationships among
    loss, parameter count, training data, and compute.
    The important engineering question is not simply how large a model can be,
    but how a fixed compute budget should be divided between model capacity and
    training tokens.

    \factvisual{
        \textbf{The Chinchilla correction.}
        Hoffmann et al.\ trained more than 400 models and found that many large
        language models were undertrained.
        For compute-optimal training, doubling parameter count should be paired
        with roughly twice as many training tokens.

        At the same training-compute budget as the 280B-parameter Gopher,
        \textbf{Chinchilla used 70B parameters and four times more data}.
        It outperformed Gopher and several other larger models across a broad
        evaluation suite.\textsuperscript{\cite{hoffmannTrainingComputeOptimal2022}}
    }{
        \begin{tikzpicture}[x=5.5mm,y=8mm]
            \node[anchor=east,font=\scriptsize\sffamily\bfseries]
                at (0,2.2) {Gopher};
            \node[anchor=east,font=\scriptsize\sffamily\bfseries]
                at (0,0.6) {Chinchilla};
            \node[anchor=south,font=\scriptsize\sffamily,
                  text=AlmanachTeal] at (2,3) {model size};
            \node[anchor=south,font=\scriptsize\sffamily,
                  text=AlmanachOchre] at (6.5,3) {training data};
            \foreach \x in {0.5,1.5,2.5,3.5}{
                \fill[AlmanachTeal!75] (\x,1.9) rectangle +(0.75,0.6);
            }
            \fill[AlmanachOchre!75] (5.3,1.9) rectangle +(0.75,0.6);
            \fill[AlmanachTeal!75] (0.5,0.3) rectangle +(0.75,0.6);
            \foreach \x in {5.3,6.3,7.3,8.3}{
                \fill[AlmanachOchre!75] (\x,0.3) rectangle +(0.75,0.6);
            }
            \node[font=\tiny,anchor=west] at (0.5,1.55)
                {$4\times$ relative parameters};
            \node[font=\tiny,anchor=west] at (5.3,1.55)
                {$1\times$ relative data};
            \node[font=\tiny,anchor=west] at (0.5,-0.05)
                {$1\times$ relative parameters};
            \node[font=\tiny,anchor=west] at (5.3,-0.05)
                {$4\times$ relative data};
        \end{tikzpicture}
        \visualcaption{At equal training compute, Chinchilla reallocated the
          budget from parameter count to substantially more training data.
          Blocks show relative, not absolute, quantities.}
          {fig:chinchilla-compute-allocation}
    }

    \begin{engineernote}[What Chinchilla does not imply]
    Compute-optimal pre-training is not a universal deployment optimum.
    Data quality, inference cost, latency, memory, and the expected number of
    downstream queries can justify training a smaller model on more tokens.
    Treat the paper as a resource-allocation result under stated assumptions,
    not as a timeless constant for every architecture and dataset.
    \end{engineernote}

\end{v3topic}
\begin{v3topic}{Emergent Abilities of Large Models}
    The term is due to Wei and
    colleagues\textsuperscript{\cite{weiEmergentAbilitiesLarge2022}}; few-shot
    behaviour was first documented at scale with
    GPT-3\textsuperscript{\cite{brownLanguageModelsAre2020}}.

    \textbf{Emergent abilities} are capabilities that appear only above a certain model
    scale and are \emph{not} present in smaller models --- they cannot be predicted by
    extrapolating smaller-scale
    behaviour.\textsuperscript{\cite[][p.~107]{farrisHowLargeLanguage2025}}

    Examples observed in GPT-3 and beyond:
    \begin{itemize}
        \item \textbf{Few-shot in-context learning:} solving tasks from a handful of
              examples in the prompt with no weight updates.
        \item \textbf{Chain-of-thought reasoning:} multi-step arithmetic and logical
              deduction when prompted to ``think step by step.''
        \item \textbf{Code generation:} producing syntactically and semantically correct
              programs in dozens of languages.
        \item \textbf{Instruction following:} executing complex, multi-part natural
              language instructions without explicit training on those instruction formats.
    \end{itemize}
    These emergent properties make LLMs qualitatively different from smaller
    pre-transformer NLP
    models.\textsuperscript{\cite[][p.~114]{farrisHowLargeLanguage2025}}

\end{v3topic}

% ---- 12.4.4 Prompt Engineering --------------------------------------------------------------
\subsection{Prompt Engineering}


\begin{v3topic}{Prompt Frameworks: The Named Mnemonics}
    Alongside the \emph{techniques} below --- which are research results about
    how models behave --- a parallel body of \keyterm{prompt frameworks} has
    grown up in practice. These are mnemonics: checklists of the dimensions a
    prompt should specify, packaged as an acronym so a non-specialist can
    recall them.

\begin{table}[htbp]
\centering
\caption{Prompt frameworks in common circulation. Read down the expansions
rather than across the names: the same handful of dimensions --- role,
context, task, constraints, format, audience, evaluation --- recurs in almost
every row.}
\label{tab:nlp-prompt-frameworks}
\begin{tblr}{
  width=\textwidth,
  colspec={Q[l,0.155\textwidth] X[l] Q[l,0.215\textwidth]},
  row{1}={font=\bfseries,bg=AlmanachMist},
  column{1}={cmd=\nohyphens},
  hlines,vlines}
Name & Expansion & Suits \\
R-T-F & Role, Task, Format & The shortest useful checklist; everyday requests \\
A-P-E & Action, Purpose, Expectation & Speed over precision \\
T-A-G & Task, Action, Goal & Simple task framing \\
R-A-C-E & Role, Action, Context, Expectation & Quick structured prompts \\
C-A-R-E & Context, Action, Result, Example & When one worked example carries
  the intent \\
C-L-E-A-R\textsuperscript{\cite{loCLEARPathFramework2023}} &
  Concise, Logical, Explicit, Adaptive, Reflective &
  Not a slot list but a set of writing \emph{qualities}; the one framework
  here with peer-reviewed provenance \\
R-I-S-E-N & Role, Instructions, Steps, End goal, Narrowing &
  Multi-step procedural tasks \\
S-T-O-K-E & Situation, Task, Objective, Knowledge, Examples &
  Domain-expert tasks where accuracy needs supplied knowledge \\
C-R-I-S-P-E & Capacity/role, Insight, Statement, Personality, Experiment &
  Exploratory work; asks for several variants \\
T-C-R-E-I\textsuperscript{\cite{googlePromptingEssentials2026}} &
  Task, Context, References, Evaluate, Iterate &
  Includes the two steps most people skip: check the output, then revise the
  prompt \\
C-O-S-T-A-R\textsuperscript{\cite{teoHowWonSingapore2023}} &
  Context, Objective, Style, Tone, Audience, Response format &
  Audience-facing writing, where tone and register matter \\
C-R-E-A-T-E & Character, Request, Examples, Adjustments, Type, Extras &
  Creative work matched to a reference style \\
B-R-I-D-G-E & Background, Request, Inputs, Deliverable, Guardrails,
  Evaluation & Work with supplied data, hard constraints, and a self-check \\
5W3H & What, Why, Who, When, Where, How-to-do, How-much, How-feel &
  The most exhaustive; derived from the journalistic
  checklist\textsuperscript{\cite{penggangStructuredIntentProtocol2026}} \\
\end{tblr}
\end{table}

\begin{plainlanguage}
None of these is an algorithm. They are packing lists. The reason a prompt
fails is almost never that the writer chose the wrong acronym --- it is that
they left out the audience, or the format, or the constraint that the answer
must fit on one page. An acronym is simply a way of not forgetting.
\end{plainlanguage}

\begin{workedexample}[One task, three frameworks, and what actually changed]
\emph{The bare request:} ``Write something about our new API rate limits.''

\emph{R-T-F.}
\textsf{Role:} you are a developer-relations engineer.
\textsf{Task:} explain the new rate limits and what a caller should do when
they hit one.
\textsf{Format:} 200 words, then a three-row table of limit tiers.

\emph{C-O-S-T-A-R} adds three slots R-T-F does not have:
\textsf{Context:} limits change on 1 October; existing integrations will break
without action.
\textsf{Objective:} get integrators to add retry-with-backoff before that date.
\textsf{Style:} technical documentation.
\textsf{Tone:} matter-of-fact, not apologetic.
\textsf{Audience:} third-party developers who did not ask for this change.
\textsf{Response:} 200 words plus a table.

\emph{T-C-R-E-I} adds two \emph{process} steps rather than content slots:
\textsf{References:} last quarter's deprecation notice, as a style exemplar.
\textsf{Evaluate:} check every stated number against the specification.
\textsf{Iterate:} revise the prompt, not just the output.

\emph{What each addition bought.} R-T-F fixed the shape. CO-STAR fixed the
\emph{stance} --- ``matter-of-fact, not apologetic'' and ``developers who did
not ask for this'' change the text substantially, and neither is expressible
in R-T-F. TCREI fixed the \emph{workflow}: the reference sets the house style
without describing it, and the evaluate step catches the invented rate number
that no amount of prompt structure prevents.

\emph{The point.} The gain came from the dimensions added, not from the
mnemonic. Had you written the same six considerations as unnumbered prose,
you would have got the same result.
\end{workedexample}

\begin{engineernote}[What the evidence actually says]
A controlled comparison across three frontier models, three languages, and
$3{,}240$ generated outputs found that 5W3H, CO-STAR, and RISEN reached
\emph{statistically comparable} goal-alignment scores ($4.93$, $4.98$, $4.98$
on a five-point scale), leading the authors to conclude that
\emph{dimensional decomposition itself} is the active ingredient rather than
any particular
framework\textsuperscript{\cite{penggangStructuredIntentProtocol2026}}.

Two further findings are more useful than the ranking. First, structured
prompting cut cross-language variance dramatically ($\sigma$ from $0.470$ to
about $0.020$): the same intent produced far more consistent output regardless
of the language it was written in. Second --- the \keyterm{weak-model
compensation effect} --- the weakest baseline model gained $+1.006$ points from
structure while the strongest gained only $+0.217$.

\emph{Read that second result carefully, because it is the practical one.}
Elaborate prompt scaffolding buys most where the model is weakest, and its
value shrinks as models improve. A framework that was essential on a 2023
model may be redundant on a current one, and the way to find out is to test
the plain request first.

\emph{And read the study's limits.} It is a single-author preprint, not peer
reviewed; its judge is itself a language model; and the authors explicitly
note that the near-identical scores may reflect a ceiling in the metric rather
than genuine equivalence. It is the best available evidence on this question,
which is not the same as strong evidence.
\end{engineernote}

\begin{v3topic}{How Much Does Prompting Actually Buy?}
    Prompting advice is usually asserted rather than measured. The published
    numbers that do exist span three orders of magnitude, and the spread is the
    finding: what you gain depends far more on the \emph{task} and the
    \emph{model} than on the framework.

\begin{table}[htbp]
\centering
\caption{Measured effects of prompting interventions. Read the last two rows
against the first: the largest published gains come from reasoning
\emph{techniques} on hard tasks, not from prompt \emph{frameworks}, and one
study finds structure that helps in isolation collapsing in a realistic
prompt.}
\label{tab:nlp-prompting-effectiveness}
\begin{tblr}{
  width=\textwidth,
  colspec={Q[l,0.21\textwidth] X[l] Q[l,0.20\textwidth]},
  row{1}={font=\bfseries,bg=AlmanachMist},
  hlines,vlines}
Intervention & Measured effect & Setting \\
Chain-of-thought\textsuperscript{\cite{weiChainofThoughtPrompting2022}} &
  GSM8K solve rate $18\,\% \rightarrow 57\,\%$ ($+39$ points) &
  PaLM 540B; gain is \emph{emergent} --- small models get little or none \\
Self-consistency over CoT\textsuperscript{\cite{wangSelfConsistencyImprovesChain2023}} &
  GSM8K $+17.9$, SVAMP $+11.0$, AQuA $+12.2$, StrategyQA $+6.4$,
  ARC-challenge $+3.9$ points &
  Costs $n$ samples per answer \\
Structured clinical reasoning\textsuperscript{\cite{sonodaStructuredClinicalReasoning2024}} &
  Primary diagnosis $56.5\,\% \rightarrow 60.6\,\%$ ($p=0.042$); top-three
  $66.5\,\% \rightarrow 70.5\,\%$ ($p=0.005$) &
  322 radiology cases; a summary-only control gave $56.2\,\%$, i.e.\ no gain \\
Framework choice among CO-STAR, RISEN, 5W3H\textsuperscript{\cite{penggangStructuredIntentProtocol2026}} &
  Goal alignment $4.93$ / $4.98$ / $4.98$ of 5 --- \textbf{no meaningful
  difference} & 3{,}240 outputs, 3 models, 3 languages \\
Any structure vs.\ none, by model strength\textsuperscript{\cite{penggangStructuredIntentProtocol2026}} &
  Weakest model $+1.006$; strongest model $+0.217$ &
  The gain shrinks as the model improves \\
Any structure vs.\ none, cross-language consistency\textsuperscript{\cite{penggangStructuredIntentProtocol2026}} &
  Cross-language $\sigma$ $0.470 \rightarrow 0.020$ &
  Consistency improves far more than mean quality \\
Structure inside a busy production prompt\textsuperscript{\cite{joPromptComplexityDilutes2026}} &
  $100\,\% \rightarrow 0\text{--}30\,\%$ when competing instructions are added
  & The same structure, embedded rather than isolated \\
\end{tblr}
\end{table}

\begin{workedexample}[Reading the table as an engineer]
Four conclusions follow, and none of them is ``pick framework X''.

\emph{1. The big wins are techniques, not frameworks.} Chain-of-thought moved
GSM8K by $39$ points. The best-measured \emph{framework} comparison found a
difference of $0.05$ on a five-point scale between three named frameworks ---
statistically nothing. If you have a hard reasoning task, spend your effort on
the technique.

\emph{2. Gains are task-shaped.} $+39$ points on multi-step arithmetic;
$+4.1$ points on radiological diagnosis. Both are real; quoting the first as
representative of prompting in general is not.

\emph{3. Structure buys consistency more reliably than quality.} The
cross-language variance fell by a factor of about $24$ while mean scores moved
comparatively little. For a production system serving many users, \emph{lower
variance is often worth more than a higher mean} --- it is the difference
between a system that is usually good and one that is dependably adequate.

\emph{4. The benefit is being competed away.} The weak model gained nearly five
times what the strong one did. Extrapolate: advice written against a 2023 model
may be measuring a deficiency that no longer exists. Always test the plain
request as your baseline before adopting elaborate scaffolding.
\end{workedexample}

\begin{cautionbox}[Structure that works in isolation can fail in place]
The most useful negative result here is also the least known. A structured
reasoning format that scored $100\,\%$ on a task in isolation dropped to
between $0$ and $30\,\%$ when the same structure was placed inside a realistic
production system prompt alongside other
instructions\textsuperscript{\cite{joPromptComplexityDilutes2026}}. The
diagnosed mechanism is instructive: a competing directive demanded a
conclusion first, which inverted the reason-then-answer order the structure
depends on. The model was observed emitting its answer \emph{before} the
reasoning it was told to perform --- it retained the capability and had already
committed.

Two lessons. \textbf{Instruction order is a design variable}, not formatting.
And \textbf{benchmark your prompt where it will live}, not on its own: an
isolated A/B test of a prompt fragment can be actively misleading about its
behaviour in the full system. This is a single preprint and should be treated
as a caution to test for rather than an established magnitude --- but the
failure mode is easy to check and cheap to prevent.
\end{cautionbox}
\end{v3topic}

\begin{v3topic}{Choosing and Testing a Prompt}
\begin{cautionbox}[These acronyms are folklore, and they do not agree with each other]
Unlike the techniques in the next topic, prompt frameworks are largely
practitioner artefacts with no common standard, and their expansions drift.
R-I-S-E-N appears in circulation as \emph{Role, Instructions, Steps, End goal,
Narrowing}; as \emph{Role, Input, Steps, Expectation, Narrowing}; and as
\emph{Role, Input, Scenario, Expectation, Nuance} --- three different
checklists behind one name. C-L-E-A-R is a set of writing qualities, while
C-O-S-T-A-R is a set of content slots; they are not the same kind of object
despite being presented side by side.

Two consequences. Do not cite a framework as though it were a standard ---
state the expansion you mean. And do not let framework choice absorb effort
that belongs in evaluation: a prompt written to \emph{any} of these and then
tested against twenty fixed cases will beat a prompt written to the
theoretically best one and never measured.
\end{cautionbox}

\begin{engineernote}[When a framework stops being the right tool]
Frameworks address the \emph{writing} of a single prompt. Three situations
outgrow them:
\begin{itemize}
  \item \textbf{The prompt is now a program.} Once it is versioned, tested,
        and templated with variables, treat it as code --- with a test suite ---
        rather than as prose to be checklisted.
  \item \textbf{The context is assembled, not written.} In an agent, most of
        the context is retrieved documents, tool results, and history. That is
        \emph{context engineering} (\cref{sec:mas-engineering-llm-agents}), and
        a mnemonic about role and tone does not touch it.
  \item \textbf{The task needs a different technique, not a better prompt.}
        If the model lacks the knowledge, retrieval is the answer
        (\cref{fig:nlp-rag-pipeline}); if it lacks the behaviour, fine-tuning
        or an adapter is (\cref{ch:practical-llm-engineering}). No amount of
        structure substitutes for either.
\end{itemize}
\end{engineernote}

\begin{checkpoint}
Take \cref{tab:nlp-prompt-frameworks} and list the distinct \emph{dimensions}
that appear across all fourteen rows, merging synonyms. You should end with
roughly eight. Then write your own checklist covering those eight, and explain
why it is neither better nor worse than any published framework.
\end{checkpoint}
\end{v3topic}

\begin{v3topic}{Prompting Strategies}
    Chain-of-thought prompting is due to Wei and
    colleagues\textsuperscript{\cite{weiChainofThoughtPrompting2022}}.

    \textbf{Prompting} shapes LLM behaviour without changing model
    weights.\textsuperscript{\cite[][p.~82]{farrisHowLargeLanguage2025}}

    \begin{description}[leftmargin=0.4cm, labelindent=0pt]
        \item[\textbf{Zero-shot}] Task description only; no examples.
              Tests the model's built-in generalisation.
        \item[\textbf{Few-shot}] $k$ input--output examples in the context window
              ($k=1$--$8$ typical). The model infers the task pattern without gradient
              updates --- \emph{in-context learning}.
        \item[\textbf{Chain-of-Thought (CoT)}] Adding ``Let's think step by step'' or
              worked reasoning traces dramatically improves multi-step arithmetic, logic,
              and science problems.
        \item[\textbf{Self-consistency}] Generate multiple CoT paths; majority-vote over
              answers for improved reliability.
    \end{description}

\end{v3topic}
\begin{v3topic}{System Prompts and Prompt Templates}
    A \textbf{system prompt} is prepended text that sets the model's \emph{persona},
    output format, and behavioural
    constraints.\textsuperscript{\cite[][p.~68]{farrisHowLargeLanguage2025}}

    \begin{itemize}
        \item \emph{``You are a helpful medical assistant. Always recommend consulting a
              doctor.''} --- sets persona and safety constraint.
        \item \emph{``Respond only in JSON with keys: \{result, confidence\}.''} ---
              enforces structured output.
    \end{itemize}

    \textbf{Prompt templates} are reusable structures with variable placeholders,
    enabling systematic evaluation and deployment:
    \begin{center}
        \texttt{Translate the following \{lang\} text to English: \{text\}}
    \end{center}

    Good prompt design requires clear instructions, relevant context, and explicit output
    format
    requirements.\textsuperscript{\cite[][p.~70]{farrisHowLargeLanguage2025}}

\end{v3topic}

% ---- 12.4.5 Retrieval-Augmented Generation --------------------------------------------------
\subsection{Retrieval-Augmented Generation}


\begin{v3topic}{RAG: Motivation and Pipeline}
    \textbf{Retrieval-Augmented Generation (RAG)}\textsuperscript{\cite{lewisRetrievalAugmentedGeneration2020}} addresses two core LLM limitations:
    (1) \emph{knowledge cutoff} --- the model cannot know events after training, and
    (2) \emph{hallucination} --- the model generates plausible but factually wrong
    content.\textsuperscript{\cite[][p.~82]{farrisHowLargeLanguage2025}}

    RAG augments the LLM with an external \emph{knowledge base} retrieved at inference
    time. The pipeline:

    \begin{enumerate}
        \item \textbf{Index:} split documents into chunks; compute dense vector
              embeddings; store in a vector database.
        \item \textbf{Retrieve:} embed the user query; perform approximate
              nearest-neighbour (ANN) search; return top-$k$ chunks.
        \item \textbf{Augment:} inject retrieved chunks into the LLM prompt as context
              alongside the original query.
        \item \textbf{Generate:} the LLM produces a grounded answer conditioned on the
              retrieved context, with citations if desired.
    \end{enumerate}

    RAG can be combined with fine-tuning or used with an unmodified base model. It is
    particularly effective for enterprise knowledge bases, document QA, and keeping
    responses up-to-date.\textsuperscript{\cite[][p.~83]{farrisHowLargeLanguage2025}}

\begin{figure}[htbp]
\centering
\begin{tikzpicture}[
  node distance=6mm,
  bx/.style={draw=AlmanachTeal,fill=AlmanachMist,rounded corners=1mm,
    minimum height=9mm,minimum width=21mm,font=\sffamily\scriptsize,
    align=center},
  db/.style={bx,fill=AlmanachOchre!20,draw=AlmanachOchre},
  arr/.style={-{Stealth[length=1.8mm]},thick,draw=AlmanachInk!55},
  lbl/.style={font=\sffamily\scriptsize,color=AlmanachInk!70}]

  \node[lbl,anchor=west,font=\sffamily\scriptsize\bfseries] at (0,15mm)
    {Offline: index once};
  \node[bx] (docs) at (12mm,4mm) {documents};
  \node[bx,right=of docs] (chunk) {chunk\\\scriptsize 512 tokens};
  \node[bx,right=of chunk] (emb) {embed};
  \node[db,right=of emb] (vdb) {vector\\store};
  \foreach \a/\b in {docs/chunk,chunk/emb,emb/vdb}{\draw[arr] (\a) -- (\b);}

  \node[lbl,anchor=west,font=\sffamily\scriptsize\bfseries] at (0,-14mm)
    {Online: per query};
  \node[bx] (q) at (12mm,-25mm) {query};
  \node[bx,right=of q] (qe) {embed};
  \node[bx,right=of qe] (ann) {top-$k$\\search};
  \node[bx,right=of ann] (llm) {LLM};
  \node[bx,right=of llm] (ans) {answer};
  \foreach \a/\b in {q/qe,qe/ann,ann/llm,llm/ans}{\draw[arr] (\a) -- (\b);}
  \draw[arr,draw=AlmanachOchre] (vdb) -- (ann)
    node[midway,right,lbl,color=AlmanachOchre] {$k$ chunks};
  \draw[arr,dashed,draw=AlmanachInk!40] (q.south) -- ++(0,-6mm) -| (llm.south)
    node[near end,below,lbl] {original question carried through};
\end{tikzpicture}
\caption{The two halves of a RAG system. Indexing happens once and offline;
retrieval and generation happen per query. The two halves fail independently,
which is why they must be evaluated separately.}
\label{fig:nlp-rag-pipeline}
\end{figure}

\begin{workedexample}[Sizing a RAG pipeline before building it]
A knowledge base of $50{,}000$ tokens, chunk size $512$, overlap $64$, feeding
a model with a $4{,}096$-token context.

\emph{How many chunks?} The stride is $512 - 64 = 448$ tokens, so
\[
  \text{chunks} = \left\lceil\frac{50{,}000 - 512}{448}\right\rceil + 1
  = \lceil 110.5 \rceil + 1 = 112 .
\]
The overlap costs about $12\,\%$ extra storage and exists so that a fact
straddling a boundary still appears intact in at least one chunk.

\emph{What fits in the prompt?} Retrieving $k=5$ chunks costs
$5 \times 512 = 2{,}560$ tokens. Add a system prompt and the question, say
$300$ tokens. That leaves
\[
  4096 - 2560 - 300 = 1{,}236 \text{ tokens for the answer.}
\]

\emph{Now the trade-off is visible.} Raising $k$ to $8$ costs $4{,}096$ tokens
of retrieved text alone --- the entire window, with nothing left to answer
with. Every RAG design decision is a negotiation inside this budget, and it is
arithmetic you can do before writing any code.
\end{workedexample}

\begin{cautionbox}[Retrieval failure and generation failure look identical]
When a RAG system returns something wrong, there are two completely different
causes and the output does not distinguish them: either the right chunk was
never retrieved, or it was retrieved and the model ignored or misread it. The
first is a search problem (embedding model, chunking, $k$, reranking); the
second is a grounding problem (prompt, model, instruction to abstain).

Diagnose by checking retrieval separately: for a set of questions with known
answers, measure whether the supporting chunk appears in the top $k$ at all.
That single measurement splits your debugging in half, and it is the reason
\cref{fig:nlp-rag-pipeline} is drawn as two independent halves.
\end{cautionbox}

\end{v3topic}
\begin{v3topic}{Vector Databases and Approximate Search}
    A \keyterm{vector database} stores embeddings and answers nearest-neighbour
    queries over them. The reason it is a distinct piece of infrastructure
    rather than a table with a distance function is arithmetic: exact search
    over $N$ vectors of dimension $d$ costs $O(Nd)$ per query.

\begin{workedexample}[Why exact nearest-neighbour search stops working]
Ten million chunks embedded at $d = 1{,}536$ in \texttt{float32}.

\emph{Storage.} $10^7 \times 1536 \times 4 = 61.4\,\text{GB}$ for the vectors
alone.

\emph{One exact query.} A dot product per vector is $1{,}536$ multiply--adds,
so $10^7 \times 1536 = 1.54\times10^{10}$ operations --- roughly a second of
pure compute per query, before any other work, and it scales linearly with the
corpus.

\emph{Approximate search} trades exactness for speed. HNSW builds a layered
proximity graph and walks it greedily, reaching sub-millisecond queries at
roughly $95\text{--}99\,\%$
recall\textsuperscript{\cite{malkovEfficientRobustApproximate2020}};
IVF and product-quantisation methods partition and compress the space
instead\textsuperscript{\cite{johnsonBillionscaleSimilaritySearch2021}}.

\emph{The consequence nobody plans for.} Your retrieval is now
\emph{probabilistic}. The chunk that should have been returned may simply not
be, and no error is raised. When measuring a RAG system
(\cref{fig:nlp-rag-pipeline}), measure recall against exact search on a sample
--- otherwise a retrieval regression is indistinguishable from a generation
regression.
\end{workedexample}

\begin{engineernote}
Three practical properties decide the choice, and none of them is the
similarity metric. \textbf{Filtering}: can it combine a vector query with a
metadata predicate (tenant, date, permission) \emph{during} the search rather
than after it? Post-filtering silently returns fewer results than requested.
\textbf{Updates}: can vectors be deleted and re-indexed without a full
rebuild? \textbf{Recall measurability}: can you compare against exact search?
A store that cannot answer the third leaves you unable to debug the first.

Note also that embeddings are derived data. Changing the embedding model
invalidates every stored vector, so re-embedding the whole corpus is a
migration to plan for, not an accident to discover.
\end{engineernote}

\begin{cautionbox}[Hybrid search exists because embeddings miss exact matches]
Dense retrieval matches meaning and is weak precisely where meaning is not the
point: part numbers, error codes, surnames, statute references, rare acronyms.
An embedding places \texttt{ERR-4471} near other error codes rather than
retrieving that one. Keyword search (BM25) does the opposite --- exact on
tokens, blind to paraphrase.

Production systems generally run both and fuse the rankings, then rerank the
merged top candidates with a cross-encoder. If your RAG system fails on
identifiers, the fix is not a better embedding model.
\end{cautionbox}
\end{v3topic}

\begin{v3topic}{Graph-Based Retrieval}
    Chunk-based retrieval answers questions whose evidence sits in one place.
    It fails on questions requiring several facts to be joined --- ``which of
    our suppliers are affected by the regulation that changed last quarter?''
    --- because no single chunk contains the answer and top-$k$ retrieval has
    no notion of connecting them.

    \keyterm{Graph-based retrieval} builds an explicit structure over the
    corpus: entities as nodes, relationships as edges, extracted by a model
    during indexing. Retrieval then traverses the graph rather than ranking
    passages, and community summaries over the graph can answer corpus-level
    questions that no individual chunk
    supports\textsuperscript{\cite{edgeLocalGlobalGraph2024}}.

    The costs are real and should be stated plainly: indexing requires a model
    pass over the whole corpus to extract entities and relations, extraction
    errors become permanent structural errors, and the graph must be
    maintained as documents change. Reach for it when the questions are
    genuinely multi-hop or aggregate, not by default.

\begin{engineernote}
Knowledge graphs predate all of this and remain the right tool when relations
are known rather than inferred --- product hierarchies, organisational
structure, bills of materials, clinical ontologies. The pragmatic pattern is a
curated graph for the relations you can state authoritatively, plus vector
retrieval over the unstructured text, with the graph used to constrain or
expand what is retrieved.
\end{engineernote}
\end{v3topic}

\begin{v3topic}{Chunking Strategies}
    The retrieval step depends on \textbf{dense vector embeddings}: each text chunk is
    encoded as a high-dimensional vector by an embedding model (e.g.\ a BERT-based
    sentence encoder). Similarity is measured by cosine distance or inner product.

    \textbf{Vector databases} (FAISS, Pinecone, Weaviate, ChromaDB) provide efficient ANN
    search over millions of embeddings at low
    latency.\textsuperscript{\cite[][p.~83]{farrisHowLargeLanguage2025}}

    \textbf{Chunking strategy} matters significantly:
    \par\smallskip\noindent
    \begin{tblr}{
        width=\linewidth,
        colspec={X[l,2] X[l,4]},
        hlines, vlines,
        row{1}={font=\bfseries\small, bg=LimeGreen!20},
        row{2-Z}={font=\small},
    }
        Strategy & Description \\
        Fixed-size & Split every $N$ tokens; fast and simple; may cut mid-sentence \\
        Sentence / paragraph & Split on natural boundaries; preserves semantic coherence \\
        Semantic chunking & Embed and cluster sentences by similarity; highest quality \\
        Recursive splitting & Hierarchically split: paragraph $\to$ sentence $\to$ token \\
    \end{tblr}
\captionof{table}{Chunking strategies. Chunk size trades retrieval precision against the risk of splitting a fact across a boundary --- see the sizing example in this chapter.}
\label{tab:nlp-chunking}
    Chunk size trades off recall (larger $=$ more context per chunk) against precision
    (smaller $=$ more focused retrieval).

\end{v3topic}

% ==== Section 5: Application Scenarios =======================================================
\section{Application Scenarios}

% ---- 12.5.1 Machine Translation -------------------------------------------------------------
\subsection{Machine Translation}


\begin{v3topic}{Machine Translation}
    \textbf{Neural Machine Translation (NMT)} maps a source-language sequence to a
    target-language sequence using an encoder--decoder
    Transformer.\textsuperscript{\cite[][p.~7]{dhamaniIntroductionGenerativeAI2024}}

    Evolution: rule-based (1950s--1980s) $\to$ statistical MT (IBM models, phrase-based)
    $\to$ seq2seq + attention (Bahdanau 2014) $\to$ Transformer-based NMT (Vaswani 2017).

    \textbf{Google NMT} (Wu et al.\ 2016) was the first large-scale production
    transformer-based MT system, reducing translation errors by $\sim$60\,\% relative to
    phrase-based
    SMT.\textsuperscript{\cite[][p.~9]{dhamaniIntroductionGenerativeAI2024}}

    Key quality metric: \textbf{BLEU score} measures $n$-gram overlap between hypothesis
    and reference translation. Imperfect proxy for human judgment but standard for
    benchmarking.

\end{v3topic}
\begin{v3topic}{Sequence-to-Sequence with Attention}
    Encoder--decoder architecture for
    MT:\textsuperscript{\cite[][p.~7-8]{dhamaniIntroductionGenerativeAI2024}}
    \begin{enumerate}
        \item \textbf{Encoder:} reads source sequence $x_1,\ldots,x_S$; outputs
              contextual representations $h_1,\ldots,h_S$.
        \item \textbf{Attention:} for each target position $t$, computes a weighted sum
              over encoder states:
              $c_t = \sum_s \alpha_{ts}\,h_s$.
        \item \textbf{Decoder:} generates target tokens autoregressively conditioned
              on $c_t$ and previous outputs.
    \end{enumerate}
    The Transformer replaces recurrence entirely with \emph{self-attention}, enabling
    full parallelisation and capturing long-range dependencies more effectively than
    LSTMs.\textsuperscript{\cite[][]{vaswaniAttentionAllYou2023}}

\end{v3topic}

% ---- 12.5.2 Information Extraction and Text Understanding -----------------------------------
\subsection{Information Extraction and Text Understanding}


\begin{v3topic}{Named Entity Recognition and Sentiment Analysis}
    \textbf{Named Entity Recognition (NER):} identify and classify spans in text as
    persons, organisations, locations, dates, etc. BERT fine-tuned on CoNLL-2003 achieved
    F1\,=\,92.8, outperforming task-specific
    models.\textsuperscript{\cite[][p.~6]{devlinBERTPretrainingDeep2019}}

    NER is fundamental for knowledge-graph construction, information retrieval, and
    document classification.

    \textbf{Sentiment Analysis:} classify the opinion polarity of text (positive /
    negative / neutral). Fine-tuned transformers achieve high accuracy on product
    reviews, social media, and financial news. Fine-grained \emph{aspect-based} sentiment
    analysis identifies sentiment towards specific entities within a
    text.\textsuperscript{\cite[][p.~12]{dhamaniIntroductionGenerativeAI2024}}

\end{v3topic}
\begin{v3topic}{Text Summarisation and Question Answering}
    \textbf{Summarisation} reduces a document to its essential content:
    \begin{itemize}
        \item \textbf{Extractive:} select and concatenate key sentences from the source.
        \item \textbf{Abstractive:} generate new text that paraphrases the content;
              requires seq2seq or LLM.
    \end{itemize}

    \textbf{Question Answering (QA)}
    types:\textsuperscript{\cite[][p.~13]{dhamaniIntroductionGenerativeAI2024}}
    \begin{itemize}
        \item \textbf{Extractive QA (SQuAD):} predict answer span start/end positions
              within a given passage.
        \item \textbf{Open-book generative QA:} generate a free-text answer using
              provided context.
        \item \textbf{Closed-book generative QA:} generate answers from model weights
              only (no external context).
    \end{itemize}

\end{v3topic}

% ---- 12.5.3 Chatbots and Voice Assistants ---------------------------------------------------
\subsection{Chatbots and Voice Assistants}


\begin{v3topic}{Task-Oriented Dialogue Systems}
    \textbf{Task-oriented} systems (booking, customer service) use a structured
    pipeline:\textsuperscript{\cite[][p.~11]{dhamaniIntroductionGenerativeAI2024}}
    \begin{enumerate}
        \item \textbf{NLU:} intent detection + slot filling from user utterance.
        \item \textbf{Dialogue State Tracking:} maintain a belief state over unfilled
              slots across turns.
        \item \textbf{Policy:} select the next system action given the current state
              (e.g.\ query database, ask for clarification).
        \item \textbf{NLG:} surface the response as natural language.
    \end{enumerate}
    Frameworks: Rasa (open-source) and Microsoft LUIS provide modular pipelines for
    task-oriented
    dialogue.\textsuperscript{\cite[][p.~11]{dhamaniIntroductionGenerativeAI2024}}

\end{v3topic}
\begin{v3topic}{Open-Domain and LLM-Based Assistants}
    \textbf{Open-domain} chatbots engage in free-form conversation. GPT-based systems
    (ChatGPT, Claude, Gemini) became the dominant paradigm from
    2022.\textsuperscript{\cite[][p.~1]{farrisHowLargeLanguage2025}}

    Key properties of LLM-based assistants:
    \begin{itemize}
        \item \textbf{Multi-turn context:} the full conversation history is included in
              the context window.
        \item \textbf{RLHF alignment:} shapes tone, helpfulness, and safety.
        \item \textbf{Tool use:} modern agents can call external APIs, execute code, or
              search the web.
    \end{itemize}

    Voice assistants (Siri, Alexa, Google Assistant) add ASR and TTS modules around an
    NLU/NLG
    core.\textsuperscript{\cite[][p.~11]{dhamaniIntroductionGenerativeAI2024}}

\end{v3topic}

% ---- 12.5.4 NLP in Education ----------------------------------------------------------------
\subsection{NLP in Education}


\begin{v3topic}{NLP Applications in Education}
    NLP enables personalised and scalable educational tools:
    \begin{itemize}
        \item \textbf{Automated Short-Answer Grading (ASAG):} transformers compare
              student responses to model answers; enables instant feedback at scale.
        \item \textbf{Intelligent Tutoring Systems (ITS):} systems such as ASSISTments
              adapt question difficulty to student performance using NLP-based knowledge
              tracing.
        \item \textbf{Essay scoring:} automated evaluation of coherence, grammar, and
              argumentation.
        \item \textbf{Language learning:} conversational AI for L2 practice;
              multilingual cross-lingual grading systems.
    \end{itemize}
    Key challenge: ensuring \emph{construct validity} --- measuring the intended
    competency, not surface
    features.\textsuperscript{\cite[][p.~11]{dhamaniIntroductionGenerativeAI2024}}

\end{v3topic}
\begin{v3topic}{LLMs as Educational Assistants}
    LLMs such as GPT-4 are increasingly used as \textbf{AI tutors}: Socratic
    questioning, step-by-step explanation, and personalised hints on
    demand.\textsuperscript{\cite[][p.~2]{farrisHowLargeLanguage2025}}

    Opportunities:
    \begin{itemize}
        \item Instant, 24/7 feedback on exercises.
        \item Adaptive difficulty and content personalisation.
        \item Support for students who lack access to human tutors.
    \end{itemize}

    Risks:
    \begin{itemize}
        \item Over-reliance may reduce deep learning and independent thinking.
        \item Detection of LLM-generated student submissions is an active challenge.
        \item Hallucination can mislead learners; accuracy is not guaranteed.
    \end{itemize}

\end{v3topic}

% ---- 12.5.5 NLP for Accessibility -----------------------------------------------------------
\subsection{NLP for Accessibility}


\begin{v3topic}{Text Simplification}
    \textbf{Text simplification} transforms complex text into a linguistically simpler
    version without losing core meaning, targeting low-literacy readers, language
    learners, and people with cognitive
    disabilities.\textsuperscript{\cite[][p.~11]{dhamaniIntroductionGenerativeAI2024}}

    Approaches:
    \begin{itemize}
        \item \textbf{Lexical simplification:} replace difficult words with simpler
              synonyms (\emph{utilise} $\to$ \emph{use}).
        \item \textbf{Syntactic simplification:} split complex sentences; reduce
              embedded clauses.
        \item \textbf{Neural simplification:} seq2seq and LLM-based generation trained
              on parallel corpora (Wikipedia and Simple Wikipedia).
    \end{itemize}

\end{v3topic}
\begin{v3topic}{Captioning and Speech Accessibility}
    NLP and speech processing combine to improve accessibility for deaf and hard-of-hearing
    users:\textsuperscript{\cite[][p.~11]{dhamaniIntroductionGenerativeAI2024}}

    \begin{itemize}
        \item \textbf{Automatic captioning:} real-time ASR (Whisper) produces closed
              captions; quality degrades with accents, fast speech, and background noise.
        \item \textbf{Screen readers:} TTS engines render digital content audibly for
              visually impaired users.
        \item \textbf{Augmentative and Alternative Communication (AAC):} language models
              predict next intended words from abbreviated input to support users with
              motor impairments.
        \item \textbf{Sign language recognition:} computer-vision-based gesture
              recognition for deaf communication.
    \end{itemize}

\end{v3topic}

% ==== Section 6: Challenges in NLP ===========================================================
\section{Challenges in NLP}

% ---- 12.6.1 Data Quality and Bias -----------------------------------------------------------
\subsection{Data Quality and Bias}


\begin{v3topic}{Training Data Bias}
    The scale-and-provenance argument is set out by Bender and
    colleagues\textsuperscript{\cite{benderDangersStochasticParrots2021}}.

    LLMs trained on web-scale data absorb and amplify the biases present in
    human-generated
    text.\textsuperscript{\cite[][p.~xii]{dhamaniIntroductionGenerativeAI2024}}

    Sources of bias:
    \begin{itemize}
        \item \textbf{Representation bias:} certain groups, languages, and viewpoints
              are over- or under-represented in training corpora.
        \item \textbf{Measurement bias:} labelling conventions in annotated datasets
              reflect the cultural assumptions of annotators.
        \item \textbf{Historical bias:} text reflects past inequalities (e.g.\ gendered
              job descriptions), which the model learns and reproduces.
    \end{itemize}

    Impact: biased models produce discriminatory outputs, unfair recommendations, and
    reinforce stereotypes at
    scale.\textsuperscript{\cite[][p.~xii]{dhamaniIntroductionGenerativeAI2024}}

\end{v3topic}
\begin{v3topic}{Fairness and Debiasing}
    \textbf{Fairness} in NLP requires that model performance does not systematically
    differ across demographic groups (gender, ethnicity, age,
    disability).\textsuperscript{\cite[][p.~140]{farrisHowLargeLanguage2025}}

    Mitigation strategies:
    \begin{itemize}
        \item \textbf{Data-level:} curate balanced corpora; filter toxic or biased
              content; augment with underrepresented voices.
        \item \textbf{Model-level:} adversarial debiasing; counterfactual data
              augmentation; fairness constraints during fine-tuning.
        \item \textbf{Output-level:} post-hoc filtering; calibration; red-teaming to
              identify failure modes.
    \end{itemize}

    No single approach is sufficient; fairness requires ongoing monitoring in
    deployment.\textsuperscript{\cite[][p.~140]{farrisHowLargeLanguage2025}}

\end{v3topic}

% ---- 12.6.2 Domain and Language Adaptation --------------------------------------------------
\subsection{Domain and Language Adaptation}


\begin{v3topic}{Multilingual Models}
    \textbf{Multilingual} LLMs (mBERT, XLM-R, mT5) are pre-trained on text from 100+
    languages in a single model, enabling \emph{zero-shot cross-lingual transfer}:
    fine-tune on English labels, evaluate on other
    languages.\textsuperscript{\cite[][p.~26]{farrisHowLargeLanguage2025}}

    Challenges:
    \begin{itemize}
        \item \textbf{Curse of multilinguality:} adding more languages reduces
              per-language capacity with fixed model size.
        \item \textbf{Low-resource languages:} languages with less training data show
              substantially worse performance (e.g.\ Swahili, Welsh, many indigenous
              languages).
        \item \textbf{Tokenizer equity:} BPE vocabularies trained on English-heavy data
              over-segment morphologically rich languages into many short tokens, wasting
              context capacity.
    \end{itemize}

\end{v3topic}
\begin{v3topic}{Domain Shift and Low-Resource Adaptation}
    A model pre-trained on general web text may perform poorly on \textbf{specialised
    domains} such as clinical notes, legal contracts, or scientific
    papers.\textsuperscript{\cite[][p.~170]{raschkaBuildLargeLanguage2025}}

    \textbf{Domain shift} occurs when the input distribution at inference differs from
    the training distribution.

    Adaptation strategies:
    \begin{itemize}
        \item \textbf{Continued pre-training:} further pre-train the base model on
              in-domain unlabelled text (e.g.\ PubMed for biomedical NLP; BioMedBERT).
        \item \textbf{Domain-specific fine-tuning:} fine-tune on labelled in-domain
              examples.
        \item \textbf{Parameter-efficient fine-tuning (PEFT):} LoRA, adapters,
              prefix-tuning --- update a small fraction of parameters to reduce compute
              cost.\textsuperscript{\cite[][p.~322]{raschkaBuildLargeLanguage2025}}
    \end{itemize}

\end{v3topic}

% ---- 12.6.3 Explainability and Safety -------------------------------------------------------
\subsection{Explainability and Safety}


\begin{v3topic}{Hallucination in LLMs}
    The phenomenon is surveyed by Ji and
    colleagues\textsuperscript{\cite{jiSurveyHallucinationNatural2023}}.

    \textbf{Hallucination} refers to LLM outputs that are fluent and confident but
    factually incorrect or
    fabricated.\textsuperscript{\cite[][p.~xii]{dhamaniIntroductionGenerativeAI2024}}

    Root causes:
    \begin{itemize}
        \item Next-token prediction optimises fluency, not factual accuracy.
        \item Rare or contradictory facts in training data are poorly memorised.
        \item Models lack a grounding mechanism to verify claims against world state.
    \end{itemize}

    Mitigation:
    \begin{itemize}
        \item \textbf{RAG:} ground generation in retrieved documents.
        \item \textbf{Self-consistency:} sample multiple outputs and compare.
        \item \textbf{Factuality fine-tuning:} train on preference pairs where truthful
              responses are preferred.
        \item \textbf{Human-in-the-loop:} maintain verification for high-stakes outputs.
    \end{itemize}

\end{v3topic}
\begin{v3topic}{Alignment and RLHF}
    Instruction tuning with human feedback is described by Ouyang and
    colleagues\textsuperscript{\cite{ouyangTrainingLanguageModels2022}}.

    \textbf{Alignment} is the problem of ensuring LLM behaviour matches human values and
    intentions.\textsuperscript{\cite[][p.~73]{farrisHowLargeLanguage2025}}

    \textbf{RLHF} (Reinforcement Learning from Human Feedback) is the primary alignment
    technique used in ChatGPT, Claude, and Gemini:

    \begin{enumerate}
        \item Pre-train or SFT the LLM.
        \item Collect human \emph{preference data}: raters rank two or more model
              responses for each prompt.
        \item Train a \textbf{reward model} $R_\phi$ to predict human preferences.
        \item Fine-tune the LLM with PPO (Proximal Policy Optimisation) to maximise
              $R_\phi$ while staying close to the SFT model (KL divergence penalty).
    \end{enumerate}

    RLHF reduces harmful outputs and improves instruction following, but does not
    eliminate misuse or
    hallucination.\textsuperscript{\cite[][p.~73-74]{farrisHowLargeLanguage2025}}

\end{v3topic}

% ==== Section 7: Further Reading =============================================================
%
\section{Deep Dive: Retrieval, Tools, and Structured Generation}
\begin{v3topic}{RAG Has Two Evaluations}
Evaluate retrieval with relevance, recall, ranking, freshness, permissions, and
coverage.
Evaluate generation with grounding, citation correctness, completeness, and
task success.
A faithful answer to irrelevant evidence is still a retrieval failure.

\end{v3topic}
\begin{v3topic}{Structured Output Is a Contract}
Define a schema, validate every field, reject unknown actions, and separate
syntactic validity from semantic authorization.
Repair loops need bounded retries and telemetry; silently accepting a plausible
but invalid object moves uncertainty into downstream code.

\end{v3topic}
\begin{v3topic}{Long Context Is Not Perfect Memory}
More tokens increase cost and can dilute relevant evidence.
Compare long-context prompting with retrieval, summarisation, and explicit
state.
Test information at different positions, conflicting evidence, multilingual
inputs, and instructions embedded inside retrieved content.
\end{v3topic}

\chapterreview{Tokenise}{Represent}{Generate}{Evaluate}


\begin{cautionbox}[Fluency is not evidence]
The single most important thing to understand about a language model is that
its confidence and its correctness are produced by different mechanisms and are
almost uncorrelated. The model is trained to produce text that \emph{looks
like} its training data. A fabricated citation looks exactly like a real one,
because looking like a real one is precisely what the objective rewarded.

The practical consequences: never accept a citation, a statistic, a quotation,
an API signature, or a legal reference from a model without checking it against
the source; and design systems so that unverifiable claims are cheap to detect
--- by grounding answers in retrieved text, by requiring citations that resolve,
and by giving the model a supported way to say it does not know.
\end{cautionbox}

\begin{practicallab}[Break a tokenizer, then break a RAG system]
\textbf{Part A --- the tokenizer.} Using any tokenizer library:
\begin{enumerate}
  \item Tokenize \texttt{"strawberry"}. Count the tokens, then ask a model how
        many times the letter \texttt{r} appears. Explain the answer from the
        token split, not from ``the model is bad at counting''.
  \item Tokenize the same paragraph in English and in a language you know that
        is written in another script. Compare token counts and compute the cost
        ratio.
  \item Tokenize \texttt{1234567} and a date. Report the splits, and predict
        which arithmetic tasks will fail.
\end{enumerate}

\textbf{Part B --- retrieval.} Build the smallest possible RAG system over
twenty documents you know well.
\begin{enumerate}
  \item Write ten questions whose answers you can point to in a specific
        document.
  \item Measure \emph{retrieval} alone: for each question, is the supporting
        chunk in the top $k$? Report recall@$k$ for $k = 1, 3, 5$.
  \item Now measure \emph{generation}: given the correct chunk, does the model
        answer correctly? Report separately.
  \item Ask a question your corpus cannot answer. Does the system abstain, or
        invent? Record which.
  \item Halve the chunk size and repeat step 2. Explain the direction of the
        change.
\end{enumerate}

\textbf{Deliverable.} Two tables --- retrieval recall and generation accuracy
--- never combined into one number. \textbf{The point:} the two halves of
\cref{fig:nlp-rag-pipeline} fail independently, and a single end-to-end score
tells you nothing about which one to fix.
\end{practicallab}

\begin{checkpoint}
Your assistant answers a question about your company's refund policy with a
confident, fluent, and wrong answer. List the four places the failure could
have originated --- chunking, embedding/retrieval, prompt, or model --- and give
the single cheapest measurement that distinguishes each from the others.
\end{checkpoint}

\section{Further Reading}


\begin{v3topic}{Recommended Sources for NLP and Generative AI}
    \begin{itemize}
        \item \fullcite{dhamaniIntroductionGenerativeAI2024} --- Accessible introduction to LLMs and generative AI: how they work, applications, limitations, and responsible use. Recommended starting point.
        \item \fullcite{farrisHowLargeLanguage2025} --- Clear technical exposition of tokenisation, transformers, learning (RLHF), RAG, and the ethics of LLMs. Ideal for practitioners.
        \item \fullcite{raschkaBuildLargeLanguage2025} --- Hands-on implementation of a GPT-style LLM from scratch in PyTorch: pre-training, classification fine-tuning, instruction fine-tuning, and LoRA.
        \item \fullcite{devlinBERTPretrainingDeep2019} --- Original BERT paper: bidirectional pre-training, MLM, NSP, and fine-tuning methodology.
        \item \fullcite{vaswaniAttentionAllYou2023} --- The foundational transformer paper: multi-head self-attention, positional encoding, encoder--decoder architecture.
        \item \fullcite{bahreeGenerativeAIAction2024} --- Practical guide to deploying generative AI in enterprise applications: integration patterns, risk management, and governance.
        \item \fullcite{anticPythonNaturalLanguage2021} --- Python-based NLP cookbook: text processing, named entity recognition, sentiment analysis, and language generation with modern libraries.
        \item \fullcite{bengioNeuralProbabilisticLanguage} --- Foundational neural language model paper introducing learned word embeddings and probabilistic next-word prediction.
    \end{itemize}

\end{v3topic}
